\chapter{Linear Algebra}
\label{chap:linear-algebra}
\section{Vector Spaces and Linear Operators}
\label{sec:vector-spaces-and-linear-operators}
	
	This section introduces vector spaces, subspaces, linear independence and linear operators: the basic algebraic structures on which every subsequent chapter depends. \\
	
	Beyond their intrinsic role in linear algebra, these structures are the language in which the rest of this paper is written, as they are essential for the development of the mathematics of neural networks. \\
	
	These structures are represented computationally throughout the project by a C++ \texttt{Matrix} class implementing construction, basic operations and transposition, given in Appendix~\ref{app:matrix-class}. Using a header \texttt{.hpp} format this Matrix object will be imported and used as a structure to help implement the computational algorithms necessary to achieve this project's goal.
	
	\subsection{Vector Spaces}
		
		\indent \indent Let $V$ be a non-empty set with operation '+', called vector addition. Assume that the pair ($V,\, +$) forms an Abelian group i.e. it satisfies the following Axioms:
		
		\begin{enumerate}[label=($A_\arabic*$), leftmargin=3em]
			\setcounter{enumi}{-1}
			\item \label{ax:A0} For all $u, v \in V$, $u + v \in V$, i.e. $V$ is closed under vector addition.
			\item \label{ax:A1} For all $u, v, w \in V$, $u + (v + w) = (u + v) + w$, that is, vector addition is associative.
			\item \label{ax:A2} For all $u, v \in V$, $u + v = v + u$, i.e. vector addition is commutative.
			\item \label{ax:A3} There is a zero element $0_V \in V$ such that $v + 0_V = v$ for all $v \in V$.
			\item \label{ax:A4} Any $v \in V$ has a negative (or additive inverse) $-v \in V$ such that $-v + v = 0_V$.
		\end{enumerate}
		
		\noindent In addition, let $\mathbb{F}$ be a field (e.g. $\mathbb{F} = \mathbb{R}$) and assume that $\mathbb{F}$ acts on $V$ via a mapping $\cdot$ called scalar multiplication and is required to satisfy the following properties.
		
		\begin{enumerate}[label=($M_\arabic*$), leftmargin=3em]
			\item \label{ax:M1} For all $\mu \in \mathbb{F}$ and $u, v \in V$, $\mu \cdot (u + v) = \mu \cdot u + \mu \cdot v$, that is, scalar multiplication is distributive over vector addition.
			\item \label{ax:M2} For all $\mu, \nu \in \mathbb{F}$ and $v \in V$, $(\mu + \nu) \cdot v = \mu \cdot v + \nu \cdot v$, that is, scalar multiplication is distributive over scalar addition.
			\item \label{ax:M3} For all $\mu, \nu \in \mathbb{F}$ and $v \in V$, $\mu \cdot (\nu \cdot v) = (\mu\nu) \cdot v$, that is, scalar multiplication is associative.
			\item \label{ax:M4} For $1 \in \mathbb{F}$, $1 \cdot v = v$ for all $v \in V$.
		\end{enumerate}
		
		\begin{definition}[Vector Space]
			A vector space over a field $\mathbb{F}$ is a set $V$ along with addition and scalar multiplication on $V$, i.e. $\left(V, +, \cdot \right)$ that obeys all axioms above. Elements of $V$ are vectors and elements of $\mathbb{F}$ are scalars.
		\end{definition}
		
		\noindent The element $0_v \in V$ is called the zero-vector, while the element $0 \in \mathbb{F}$ is called the zero-scalar. Further, subtraction in a vector space is defined as $u + (- v) \; \forall \; u, v \in V$. That is, the sum of an element with an inverse element yields a subtraction.
		
		\begin{proposition}
			\label{prop:vspace-properties}
			Let $V$ be a $\mathbb{F}$-vector space
			
			\begin{enumerate}[label=\roman*.]
				\item For all $\mu \in \mathbb{F}, \, \mu \cdot 0_v = 0_v$ and for all $v \in V, \, 0 \cdot v = 0_v$
				\item If $\mu \cdot v = 0_v$ and $\mathbb{F} \ni \mu \not= 0$, then $V \ni v = 0_v$.
				\item For all $v \in V$, $(-1) \cdot v = -v$
			\end{enumerate}
			
		\end{proposition}
		
		\begin{proof} This can be proven using the axioms ($A_0$)-($A_4$) and ($M_1$)-($M_4$).
			
			\begin{enumerate}[label=\roman*.]
				\item From axioms ($A_3$) and ($M_1$), $\mu \cdot 0_V = \mu \cdot (0_V + 0_V) = \mu \cdot 0_V + \mu \cdot 0_V$, and hence, by ($A_3$), $\mu \cdot 0_V = 0_V$. Likewise, from ($M_2$), $0 \cdot v = (0 + 0) \cdot v = 0 \cdot v + 0 \cdot v$, and hence, by ($A_3$), $0 \cdot v = 0_V$.
				\item As $\mathbb{F} \ni \mu \neq 0$, there exists $\mu^{-1} \in \mathbb{F}$ with $\mu^{-1}\mu = 1$. From axioms ($M_4$) and ($M_3$), $v = 1 \cdot v = (\mu^{-1}\mu) \cdot v = \mu^{-1} \cdot (\mu \cdot v)$. Thus, if $\mu \cdot v = 0_V$ we find that $v = \mu^{-1} \cdot 0_V$, and by (i), $v = 0_V$.
				\item From (i) and axioms ($M_2$) and ($M_4$), $0_V = 0 \cdot v = (1 + (-1)) \cdot v = 1 \cdot v + (-1) \cdot v = v + (-1) \cdot v$, and hence $(-1) \cdot v = -v$.
			\end{enumerate}
			
		\end{proof}
		
		\begin{example}[Matrices]
			The set of all $\left(m \times n\right)$-matrices with entries in $\mathbb{F}$, Mat$(m,n,\mathbb{F})$, is a $\mathbb{F}$-vector space with the usual definitions of adding matrices and multiplying by constants. For $m = n$ we write Mat$(n, \mathbb{F})$.
		\end{example}
		
	\subsection{Vector Subspaces}
		
		\indent \indent A non-empty subset $U \subseteq V$ of a vector space need not be a vector space itself, since in general it is not closed under vector addition/scalar multiplication. For example, for any $u_1, u_2 \in U$ it is not necessarily true that $u_1 + u_2 \in U$.
		
		\begin{definition}[Vector Subspace]
			Let $V$ be a $\mathbb{F}$-vector space. A vector subspace of $V$ is a non-empty subset $U$ of $V$ such that $U$ is a $\mathbb{F}$-vector space under the same operations as in $V$.
		\end{definition}
		
		\noindent If $U$ is a subset of a vector space $V$, then axioms \ref{ax:A1}, \ref{ax:A2} and \ref{ax:M1}–\ref{ax:M4} automatically hold in $U$ because they hold throughout $V$.
		
		\begin{definition}[Linear Combination]
			Let $V$ be a $\mathbb{F}$-vector space. A linear combination of the vectors $v_1, \dots, v_m \in V$ is a vector of the form 
			\[
				v = \sum_{i=1}^{m} \mu_i v_i \in V \qquad \text{with} \quad \mu_1, \dots \mu_m \in \mathbb{F}
			\]
		\end{definition}
		
		\begin{proposition}[Subspace Test]
			Let $V$ be a $\mathbb{F}$-vector space and let $U$ be a non-empty subset of $V$. Then $U$ is a vector subspace of $V$ if and only if 
			\[
				\mu_1 u_1 + \mu_2 u_2 \in U \qquad \forall \: u_1, u_2 \in U \; \textnormal{and scalars} \; \mu_1, \mu_2 \in \mathbb{F}
			\]
		\end{proposition}
		
		\noindent Thus, to decide whether $U$ is a vector subspace we do not have to check all the axioms. It is enough to show that $U$ is closed under linear combinations and that it is non-empty. \newline
		\indent Non-emptiness can be shown by demonstrating that the zero-vector is contained in $U$ (note also \ref{ax:A3}).
		
	\subsection{Linear Independence}
		
		\begin{definition}[Span]
			\label{def:span}
			Let $V$ be a $\mathbb{F}$-vector space and $v_1, \dots, v_m \in V$. The span of the set $\{v_1, \dots, v_m\}$ is the set of all possible linear combinations of $v_1, \dots, v_m$, that is,
			\[
				\operatorname{span}_{\mathbb{F}} \{ v_1, \dots, v_m \} := \left\{ \sum_{i=1}^{m} \mu_i v_i \; \middle| \; \mu_1, \dots, \mu_m \in \mathbb{F} \right\}
			\]
			If $V = \operatorname{span}_{\mathbb{F}}\{v_1, \dots, v_m\}$, then $\{v_1, \dots, v_m\}$ is said to be a spanning set for the vector space $V$.
		\end{definition}
		
		\noindent We also note that the zero-vector $0 \in V$ is always contained in the span of any vectors $v_1, \dots, v_m \in V$ since by Proposition~\ref{prop:vspace-properties}, we have $0 = 0v_1 + \dots + 0v_m \in \operatorname{span}_{\mathbb{F}} \{ v_1, \dots, v_m \}$
		
		\begin{example}[Span in $\mathbb{R}^3$]
			In $\mathbb{R}^3$,
			\[
				\operatorname{span}_{\mathbb{R}}\left\{ (1,0,0)^T, (0,1,0)^T \right\} = \left\{ (x,y,0)^T \in \mathbb{R}^3 \;\middle|\; x, y \in \mathbb{R} \right\}
			\]
			describes the $(x,y)$-plane, which is a special instance of a two-dimensional plane in $\mathbb{R}^3$ through the ``origin'' $0 \in \mathbb{R}^3$.
		\end{example}
		
		\noindent Now, let $V$ be a $\mathbb{F}$-vector space and let $v_1, \dots, v_m \in V$.
		
		\begin{definition}[Linear Dependence]
			\label{def:lin-dep}
			The vectors $v_1, \dots, v_m$ are said to be linearly dependent if and only if at least one of the vectors is a linear combination of the others.
		\end{definition} 
		
		\begin{definition}[Linear Independence]
			\label{def:lin-indep}
			The vectors $v_1, \dots, v_m$ are said to be linearly independent if and only if they are not linearly dependent.
		\end{definition}
		
		\begin{remark} 
			If a set $\{v_1, \dots, v_m\}$ contains the zero-vector $0_V \in V$, then the set is linearly dependent since $\mu \cdot 0_V = 0_V$ for all $\mu \in \mathbb{F}$ by Proposition~\ref{prop:vspace-properties}.
		\end{remark}
		
		\begin{proposition}[Extension Lemma]
			\label{prop:extension-lemma}
			Let $\{v_1, \dots, v_m\}$ be a linearly independent set of vectors in a $\mathbb{F}$-vector space $V$. Then, for any $v \in V$,
			$v \in \operatorname{span}_{\mathbb{F}}\{v_1, \dots, v_m\} \iff \{v, v_1, \dots, v_m\}$ is a linearly dependent set of vectors.
		\end{proposition}
		
		\begin{proof}
			($\Rightarrow$) Suppose $v \in \operatorname{span}_{\mathbb{F}}\{v_1, \dots, v_m\}$. Then $v = \sum_{i=1}^{m} \mu_i v_i $. By Definition~\ref{def:lin-dep}, $v$ is linearly dependent.
			
			($\Leftarrow$) Suppose $\{v, v_1, \dots, v_m\}$ is linearly dependent.
			
			\textit{Case 1:} $v$ is a linear combination of $v_1, \dots, v_m$. Then by the definition of span, $v \in \operatorname{span}_{\mathbb{F}}\{v_1,\dots,v_m\}$.
			
			\textit{Case 2:} some $v_i$ is a linear combination of the other vectors in $\{v, v_1, \dots, v_m\}$, i.e. $v_i = av + \sum_{j \neq i} b_j v_j$. If the coefficient on $v$ were zero, then $v_i$ becomes just a combination of the other $v_j$'s, contradicting linear independence. Hence that coefficient is non-zero, and solving for $v$ gives $v = a^{-1}v_i - \sum_{j\neq i} a^{-1}b_j v_j$, hence $v \in \operatorname{span}_{\mathbb{F}}\{v_1,\dots,v_m\}$.
		\end{proof}
		
		\noindent By contraposition, we now have that $V \ni v \notin \operatorname{span}_{\mathbb{F}}\{v_1, \dots, v_m\}$ if and only if $\{v, v_1, \dots, v_m\}$ forms a linearly independent set of vectors.
		
		\begin{definition}[Dimension]
			\label{def:dimension}
			A $\mathbb{F}$-vector space $V$ is said to be $n$-dimensional if and only if there is a linearly independent set of $n$ vectors and every set consisting of $n+1$ vectors is linearly dependent. $\dim_{\mathbb{F}}(V) = n$ expresses the dimension of $V$. Furthermore, a vector space is said to be of infinite dimension if and only if there is a linearly independent set of vectors of arbitrarily large order.
		\end{definition}
		
	\subsection{Bases of Vector Spaces}
		
		\indent \indent We now seek to represent any vector $v \in V$ as a linear combination of a "minimal" set of vectors, also called a basis.
		
		\begin{definition}[Basis]
			\label{def:basis}
			Let $V$ be an n-dimensional $\mathbb{F}$-vector space. Any linear independent set of n vectors in $V$ is called a \textit{basis} of $V$.
		\end{definition}
		
		\begin{corollary}[Basis Extension]
			\label{cor:basis-extension}
			Let $V$ be an $n$-dimensional $\mathbb{F}$-vector space, and let $\{v_1,\dots,v_k\}$ be a linearly independent set in $V$ with $k \le n$. Then there exist vectors $v_{k+1},\dots,v_n \in V$ such that $\{v_1,\dots,v_k,v_{k+1},\dots,v_n\}$ is a basis of $V$.
		\end{corollary}
		
		\begin{proof}
			We proceed by induction on $j$, for $j=k,k+1,\dots,n$, constructing at each stage a linearly independent set $\{v_1,\dots,v_j\}$.
			
			\textit{Base case ($j=k$).} The set $\{v_1,\dots,v_k\}$ is linearly independent by hypothesis.
			
			\textit{Inductive step.} Suppose $\{v_1,\dots,v_j\}$ is linearly independent for some $k \le j < n$. If $\{v_1,\dots,v_j\}$ spans $V$, then by Proposition~\ref{prop:v-basis} it is a basis, and since $\dim_{\mathbb{F}}(V)=n$, this forces $j=n$ — contradicting $j<n$. So $\{v_1,\dots,v_j\}$ does not span $V$, and there exists $v_{j+1} \in V \setminus \operatorname{span}_{\mathbb{F}}\{v_1,\dots,v_j\}$. By the contrapositive of Proposition~\ref{prop:extension-lemma}, $\{v_1,\dots,v_j,v_{j+1}\}$ is linearly independent, completing the inductive step.
			
			By induction, $\{v_1,\dots,v_n\}$ is linearly independent. Since $\dim_{\mathbb{F}}(V)=n$, this set of $n$ linearly independent vectors is a basis of $V$ by Definition~\ref{def:basis}. This completes the induction.
		\end{proof}
		
		\begin{proposition}
			\label{prop:v-basis}
			Let $V$ be a $\mathbb{F}$-vector space. Any linearly independent spanning set for $V$ is a basis of $V$ and vice versa.
		\end{proposition}
		
		\begin{proof}
			($\Rightarrow$) Let $\{v_1, \dots, v_n\}$ be a spanning set for $V$. Hence, any $v \in V$ lies in the span of $\{v_1, \dots, v_n\}$, that is, $v \in \operatorname{span}_{\mathbb{F}}\{v_1, \dots, v_n\}$. If, in addition, $\{v_1, \dots, v_n\}$ is linearly independent, then Proposition~\ref{prop:extension-lemma} implies that $\{v, v_1, \dots, v_n\}$ is a linearly dependent set. Thus, $\dim_{\mathbb{F}}(V) = n$ and so $\{v_1, \dots, v_n\}$ is a basis.
			
			($\Leftarrow$) Conversely, let $n = \dim_{\mathbb{F}}(V)$ and let $\{v_1, \dots, v_n\}$ be a basis of $V$. Then, for any $v \in V$, the set $\{v, v_1, \dots, v_n\}$ is linearly dependent by the definition of a basis. Finally, Proposition~\ref{prop:extension-lemma} implies that $v \in \operatorname{span}_{\mathbb{F}}\{v_1, \dots, v_n\}$, which concludes the proof.
		\end{proof}
		
		\begin{example}[Standard Basis for Mat$(m,n,\mathbb{F})$]
			The vector space Mat$(m,n,\mathbb{F})$ is spanned by the matrices $E_{ij}$ having $1$ in the $(i,j)$ position and zero elsewhere. The set of all such matrices is easily shown to be linearly independent. It is the standard basis for Mat$(m,n,\mathbb{F})$ and $\dim_{\mathbb{F}}(\text{Mat}(m,n,\mathbb{F})) = mn$.
		\end{example}
		
		\begin{theorem}[Uniqueness of Representation]
			\label{thm:unique-repr}
			Let $V$ be a $\mathbb{F}$-vector space equipped with a basis. Any vector $v \in V$ can be represented uniquely as a linear combination of the given basis vectors with coefficients in $\mathbb{F}$.
		\end{theorem}
		
		\begin{proof}
			Let $n = \dim_{\mathbb{F}}(V)$ and let $\{v_1, \dots, v_n\}$ be a basis of $V$. By virtue of Proposition~\ref{prop:v-basis}, $V = \operatorname{span}_{\mathbb{F}}\{v_1, \dots, v_n\}$, and thus any $v \in V$ can be written as $v = \sum_{i=1}^{n} \mu_i v_i$ for some $\mu_1, \dots, \mu_n \in \mathbb{F}$. To show uniqueness, let $v = \sum_{i=1}^{n} \mu_i v_i$ and $v = \sum_{i=1}^{n} \nu_i v_i$ be any two representations of $v \in V$ in terms of the basis vectors. Then, equating both expressions, we find that $\sum_{i=1}^{n} (\mu_i - \nu_i) v_i = 0$. Since the basis vectors are linearly independent, all the coefficients must vanish, that is, $\mu_i = \nu_i$ for all $i \in \{1, \dots, n\}$.
		\end{proof}
		
	\subsection{Linear Operators}
		
		\begin{definition}[Linear Operator]
			Let $V$ and $W$ be two $\mathbb{F}$-vector spaces. A linear operator is a mapping $T : V \to W$ with $T(\mu u + \nu v) = \mu T(u) + \nu T(v)$ for all $\mu, \nu \in \mathbb{F}$ and $u, v \in V$. Here, $V$ is called the domain of $T$, $W$ the codomain, and $T(v)$ the image of $v \in V$ under $T$.
		\end{definition}
		
		\noindent Linear operators are also called linear transformations, linear maps or homomorphisms. The set of all linear operators from $V$ to $W$ is denoted by $\operatorname{Hom}_{\mathbb{F}}(V,W)$. For any $T_1, T_2 \in \operatorname{Hom}_{\mathbb{F}}(V,W)$, we define 
		\[
			(\mu_1 T_1 + \mu_2 T_2)(v) := \mu_1 T_1(v) + \mu_2 T_2(v) \quad \text{for all} \: v \in V \: \text{and} \: \mu_1, \mu_2 \in \mathbb{F}
		\]
		
		\begin{example}[Transposition as a Linear Operator]
			Let $T : \text{Mat}(m,n,\mathbb{F}) \to \text{Mat}(n,m,\mathbb{F})$ be given by transposition $T : A \mapsto A^T$. For any $A, B \in \text{Mat}(m,n,\mathbb{F})$ and $a, b \in \mathbb{F}$, we have
			\(
				T(aA + bB) = (aA + bB)^T = aA^T + bB^T = aT(A) + bT(B),
			\)
			so $T$ is a linear operator.
		\end{example}
		
		\begin{proposition}[Matrix Representation of Linear Maps]
			\label{prop:matrix-representation}
			Let $T: \mathbb{F}^n \to \mathbb{F}^m$ be a linear map. Then there exists a unique matrix $A \in \mathrm{Mat}(m, n, \mathbb{F})$ such that $T(x) = Ax$ for all $x \in \mathbb{F}^n$, where the $j$-th column of $A$ is $T(e_j)$ and $\{e_1, \dots, e_n\}$ is the standard basis of $\mathbb{F}^n$.
		\end{proposition}
		
		\begin{proof}
			For any $x \in \mathbb{F}^n$, we write $x = \sum_{j=1}^n x_j e_j$. By the linearity of $T$, we have $T(x) = T(\sum_{j=1}^n x_j e_j) = \sum_{j=1}^n x_j T(e_j)$. Since matrix-vector multiplication $Ax$ is defined as the linear combination of the columns of $A$ weighted by the components of $x$, and the $j$-th column of $A$ is defined as $T(e_j)$, it follows that $T(x) = Ax$. \newline
			\indent For uniqueness, suppose $A'\in\mathrm{Mat}(m,n,\mathbb{F})$ also satisfies $T(x)=A'x$ for all $x\in\mathbb{F}^n$. For each $j$, taking $x=e_j$ gives $Ae_j = T(e_j) = A'e_j$; since $Me_j$ is the $j$-th column of $M$, $A$ and $A'$ have the same columns, so $A=A'$.
		\end{proof}
		
		\begin{corollary}[Composition and Matrix Multiplication]
			\label{cor:composition-multiplication}
			Let $T: \mathbb{F}^n \to \mathbb{F}^m$ have matrix representation $A$ and $S: \mathbb{F}^m \to \mathbb{F}^k$ have matrix representation $B$. Then the composition $S \circ T: \mathbb{F}^n \to \mathbb{F}^k$ has the matrix representation $BA$.
		\end{corollary}
		
		\begin{proof}
			For any $x \in \mathbb{F}^n$, applying Proposition~\ref{prop:matrix-representation} directly yields $(S \circ T)(x) = S(T(x)) = S(Ax)$. Applying the proposition a second time to the action of $S$ on the vector $Ax \in \mathbb{F}^m$ gives $S(Ax) = B(Ax)$. By the associativity of matrix multiplication, $B(Ax) = (BA)x$. Thus, the composition $(S \circ T)$ is represented by the matrix product $BA$.
		\end{proof}
		
		\begin{definition}[Kernel and Image]
			Let $T \in \operatorname{Hom}_{F}(V,W)$. The \textbf{kernel} of $T$ is the set $\ker(T) := \left\{ v \in V | T(v) = 0 \right\} \subseteq V$. The \textbf{image} of T is the set $\operatorname{im}(T) := \left\{ T(v) \in W | v \in V \right\} \subseteq W$.
		\end{definition}
		
		\begin{remark}
			Let $T : V \to W$ be a linear map between $\mathbb{F}$-vector spaces. Then:
			\begin{enumerate}[label=\roman*.]
				\item $\ker(T)=\{v\in V : T(v)=0_W\}$ is a subspace of $V$.
				\item $\operatorname{im}(T)=\{T(v):v\in V\}\subseteq W$ is a subspace of $W$.
			\end{enumerate}
		\end{remark}
		
		\begin{theorem}[Rank-Nullity Theorem]
			\label{thm:rank-nullity}
			For $T \in \operatorname{Hom}_{\mathbb{F}}(V,W)$, we have
			\[
				\dim_{\mathbb{F}}(\operatorname{im}(T)) + \dim_{\mathbb{F}}(\ker(T)) = \dim_{\mathbb{F}}(V)
			\]
		\end{theorem}
		
		\begin{proof}
			Set $n := \dim_{\mathbb{F}}(V)$ and $d := \dim_{\mathbb{F}}(\ker(T))$. Let $\{v_1,\dots,v_d\}$ be a basis of $\ker(T)$. Extend this to a basis (Corollary~\ref{cor:basis-extension}) $\{v_1,\dots,v_d,u_{d+1},\dots,u_n\}$ of $V$.
			
			We show that $\{T(u_{d+1}),\dots,T(u_n)\}$ is a basis of $\operatorname{im}(T)$. This suffices, since it gives
			\[
				\dim_{\mathbb{F}}(\operatorname{im}(T)) = n-d = \dim_{\mathbb{F}}(V) - \dim_{\mathbb{F}}(\ker(T))
			\]
			
			\textbf{Linear independence.} Suppose
			\[
				0 = \sum_{i=1}^{n-d}\mu_i\, T(u_{d+i}) = T\left(\sum_{i=1}^{n-d}\mu_i u_{d+i}\right)
			\]
			Then $\sum_{i=1}^{n-d}\mu_i u_{d+i} \in \ker(T)$, so it can be written in terms of the basis $\{v_1,\dots,v_d\}$:
			\[
				\sum_{i=1}^{n-d}\mu_i u_{d+i} = \sum_{j=1}^{d}\nu_j v_j
				\quad\Longleftrightarrow\quad
				\sum_{i=1}^{n-d}\mu_i u_{d+i} - \sum_{j=1}^{d}\nu_j v_j = 0
			\]
			Since $\{v_1,\dots,v_d,u_{d+1},\dots,u_n\}$ is a basis of $V$, these vectors are linearly independent. Hence all coefficients $\mu_i$ and $\nu_j$ vanish. In particular $\mu_i=0$ for all $i$, so $\{T(u_{d+1}),\dots,T(u_n)\}$ is linearly independent.
			
			\textbf{Spanning.} Let $v\in V$. Write
			\[
				v = \sum_{j=1}^{d}\alpha_j v_j + \sum_{i=1}^{n-d}\beta_i u_{d+i}
			\]
			Applying $T$ and using $T(v_j)=0$ for all $j$ (since $v_j\in\ker(T)$):
			\[
				T(v) = \sum_{i=1}^{n-d}\beta_i\, T(u_{d+i})
			\]
			So every element of $\operatorname{im}(T)$ is a linear combination of $\{T(u_{d+1}),\dots,T(u_n)\}$; these vectors span $\operatorname{im}(T)$.
			
			Hence $\{T(u_{d+1}),\dots,T(u_n)\}$ is a basis of $\operatorname{im}(T)$, completing the proof.
		\end{proof}
	
\section{The LU Decomposition}
\label{sec:lu-decomposition}
	
	This section develops the LU decomposition, a factorisation of a matrix into lower and upper-triangular components, used to solve linear systems $Ax = b$ efficiently once computed. \\
	
	Beyond its role in solving linear systems directly, this factorisation demonstrates a decompose once, then solve cheaply and repeatedly pattern. This is the same pattern that motivates the QR decomposition. \\
	
	The algorithm is implemented from scratch in C++ as \texttt{lu\_solver.cpp} (Appendix~\ref{app:lu-solver}), including partial pivoting for numerical stability.
	
	\subsection{Gaussian Elimination}
		
		\indent \indent Consider a system of $m$ linear equations in $n$ unknowns $x_1, x_2, \dots, x_n$, written in matrix form $Ax = b$, or explicitly, as:
	
		\begin{equation*}
			\begin{bmatrix}
				A_{11} & A_{12} & \dots & A_{1n} \\
				A_{21} & A_{22} & \dots & A_{2n} \\
				\vdots & \vdots & \ddots & \vdots \\
				A_{m1} & A_{m2} & \dots & A_{mn} \\
			\end{bmatrix}
			\begin{bmatrix}
				x_1 \\
				x_2 \\
				\vdots \\
				x_n \\
			\end{bmatrix}
			=
			\begin{bmatrix}
				b_1 \\
				b_2 \\
				\vdots \\
				b_n \\
			\end{bmatrix}
		\end{equation*}
		
		\noindent The key idea of Gaussian elimination is to transform a system of linear equations into a simpler system, while keeping fixed the set of solutions. This leads to the following definition.
		
		\begin{definition}[Elementary Row Operations]
			The following operations are called elementary row operations on a matrix $A$:
			\begin{enumerate}[label=\roman*.]
				\item Interchanging two rows.
				\item Multiplying a row by a non-zero scalar.
				\item Adding a multiple of one row to another row.
			\end{enumerate}
			A matrix $A$ is row-equivalent to another matrix $A'$ if and only if it can be transformed to $A'$ via elementary row operations. 
		\end{definition}
		
		\noindent Gaussian elimination is the repeated application of elementary row operations, mainly used to zero out entries below each pivot.
		
		\begin{definition}[Pivot]
			The leading/first non-zero entry in a row.
		\end{definition}
		
	\subsection{The LU Factorisation}
		
		\noindent It is often convenient to decompose a matrix $A$ into a product of two other matrices. This is called a $LU$ decomposition. We restrict our attention to square, invertible matrices $A \in \text{Mat}(n, \mathbb{F})$, since this is the setting in which the system $Ax = b$ has a unique solution.
		
		\begin{definition}[LU Decomposition]
			\label{def:lu-decomposition}
			Let $A \in \text{Mat}(n, \mathbb{F})$. An LU decomposition of $A$ is a factorisation
			\[
				A = LU
			\]
			where $L \in \text{Mat}(n, \mathbb{F})$ is unit lower triangular (that is, lower triangular with $1$'s on the diagonal) and $U \in \text{Mat}(n, \mathbb{F})$ is upper triangular.
		\end{definition}
		
		\begin{definition}[Leading Principal Minor]
			Let $A \in \text{Mat}(n, \mathbb{F})$. For $k = 1, \dots, n$, the $k$-th leading principal minor of $A$, denoted $\Delta_k$, is the determinant of the $k \times k$ submatrix formed by the first $k$ rows and first $k$ columns of $A$, that is,
			\[
				\Delta_k := \det \begin{bmatrix} 
					A_{11} & \dots & A_{1k} \\ 
					\vdots & \ddots & \vdots \\ 
					A_{k1} & \dots & A_{kk} 
					\end{bmatrix}
			\]
		\end{definition}
		
		\begin{theorem}[Existence and Uniqueness of the LU Decomposition]
			\label{thm:lu-existence}
			Let $A \in \text{Mat}(n, \mathbb{F})$ be invertible. Then $A$ admits a unique LU decomposition if and only if every leading principal minor of $A$ is non-zero, that is, $\Delta_k \neq 0$ for all $k = 1, \dots, n$.
		\end{theorem}
		
		\begin{proof}
			A full treatment of this proof is given in \cite{GolubVanLoan2013}.
		\end{proof}
		
		\begin{example}[A Zero Pivot]
			\label{ex:zero-pivot}
			Consider the matrix
			\[
				A = \begin{bmatrix} 0 & 1 \\ 1 & 0 \end{bmatrix}
			\]
			This matrix is invertible, since $\det(A) = -1 \neq 0$. However, its leading principal minors are
			\[
				\Delta_1 = A_{11} = 0, \qquad \Delta_2 = \det(A) = -1,
			\]
			and since $\Delta_1 = 0$, by Theorem~\ref{thm:lu-existence} there does not exist a unique LU decomposition for $A$.
		\end{example}
		
		\begin{definition}[Permutation Matrix]
			\label{def:permutation-matrix}
			A matrix $P \in \mathrm{Mat}(n, \mathbb{F})$ is called a \emph{permutation matrix} if it is obtained from the identity matrix $\mathbbm{1}$ by permuting its rows. That is, there exists a permutation $\sigma \in S_n$ such that
			\[
				P_{ij} = \begin{cases} 
							1 & \text{if } j = \sigma(i) \\ 
							0 & \text{otherwise}
						\end{cases}
			\]
			Equivalently, each row and each column of $P$ contains exactly one entry equal to $1$, with all other entries $0$.
		\end{definition}
		
		\noindent To extend the definition of a $LU$ decomposition to any invertible square matrix, we use the permutation matrix $P$ that acts on $A$. This gives the following extension of Theorem~\ref{thm:lu-existence}.
		
		\begin{theorem}[Existence of the $PA=LU$ Decomposition]
			\label{thm:pa-lu-existence}
			Let $A \in \mathrm{Mat}(n, \mathbb{F})$ be invertible. Then there exists a permutation matrix $P \in \mathrm{Mat}(n, \mathbb{F})$ such that $PA$ has all leading principal minors non-zero, and hence, by Theorem~\ref{thm:lu-existence}, a unique decomposition
			\[
				PA = LU,
			\]
			where $L$ is unit lower triangular and $U$ is upper triangular.
		\end{theorem}
		\begin{proof}
			A full treatment of this proof is given in \cite{GolubVanLoan2013}.
		\end{proof}
		
	\subsection{Solving the Linear System}
		We now seek to solve the system of linear equations $Ax=b$ via a $LU$ decomposition.
		
		\begin{proposition}[Solving $Ax=b$ via $PA=LU$]
			\label{prop:solve-pa-lu}
			Let $A \in \mathrm{Mat}(n,\mathbb{F})$ be invertible, let $P$ be a permutation matrix and $L,U$ unit-lower and upper triangular matrices respectively such that $PA = LU$ (Theorem~\ref{thm:pa-lu-existence}), and let $b \in \mathbb{F}^n$. Then the system
			\[
				Ax = b
			\]
			has a unique solution $x$, obtained by first solving $Ly = Pb$ for $y$ by forward substitution, then solving $Ux = y$ for $x$ by back substitution.
		\end{proposition}
		
		\begin{proof}
			$Ax=b$ is equivalent to
			\[
				PAx = Pb
			\]
			Substituting $PA = LU$,
			\[
				LUx = Pb
			\]
			Set $y := Ux$. Then the system splits into two triangular systems:
			\[
				Ly = Pb \qquad Ux = y
			\]
			Since $L$ is unit lower triangular, $L_{kk}=1\neq 0$ for all $k$, so $Ly=Pb$ has a unique solution obtained by forward substitution:
			\[
				y_k = (Pb)_k - \sum_{j=1}^{k-1} L_{kj} y_j \qquad k = 1,\dots,n
			\]
			Since $A$ is invertible and $P$ is invertible, $U = L^{-1}PA$ is invertible, hence $U_{kk} \neq 0$ for all $k$ (an upper triangular matrix is invertible if and only if all diagonal entries are non-zero). So $Ux=y$ has a unique solution obtained by back substitution:
			\[
				x_k = \frac{1}{U_{kk}}\left(y_k - \sum_{j=k+1}^{n} U_{kj}x_j\right) \qquad k = n,n-1,\dots,1
			\]
			Existence and uniqueness of $y$ and then $x$ give existence and uniqueness of a solution to $Ax=b$.
		\end{proof}
		
		\begin{remark}[The $PA=LU$ Solution Method]
			\label{rem:pa-lu-method}
			Theorem~\ref{thm:pa-lu-existence} and Proposition~\ref{prop:solve-pa-lu} combine into a single procedure for solving $Ax=b$:
			
			\begin{enumerate}[label=\roman*.]
				\item \textbf{Factor once:} compute $P$, $L$, $U$ such that $PA = LU$ (Theorem~\ref{thm:pa-lu-existence}). This step costs $O(n^3)$ operations.
				\item \textbf{Solve twice:} for any $b \in \mathbb{F}^n$, solve $Ly = Pb$ by forward substitution, then $Ux = y$ by back substitution (Proposition~\ref{prop:solve-pa-lu}). Each substitution costs $O(n^2)$.
			\end{enumerate}
			
			The point of separating factorisation from substitution is that the $O(n^3)$ cost is paid only once per matrix $A$, while each additional right-hand side $b$ costs only $O(n^2)$. This is the practical motivation for the $PA=LU$ approach over recomputing a full solve (e.g. via Cramer's rule or a fresh elimination) for every new $b$.
		\end{remark}
		
		\begin{example}[Solving $Ax=b$ via $PA=LU$]
			\label{ex:solve-pa-lu-nontrivial}
			Let
			\[
				A = \begin{bmatrix} 1 & 1 \\ 2 & 3 \end{bmatrix} \qquad
				b = \begin{bmatrix} 2 \\ 5 \end{bmatrix}
			\]
			Here $\Delta_1 = 1 \neq 0$, so $A$ admits an LU decomposition directly by Theorem~\ref{thm:lu-existence}, with $P = \mathbbm{1}$. But applying partial pivoting (Definition~\ref{def:partial-pivoting}) at step $k=1$: the active column is $(A_{11},A_{21})=(1,2)$, and $|A_{21}|=2 > |A_{11}|=1$, so the pivot row is $r=2$, and rows $1$ and $2$ are swapped before elimination proceeds.
			
			\textbf{Permutation.} Swapping rows $1$ and $2$ gives $\mathrm{perm} = [2,1]$ (in the $1$-indexed convention of \S1.2) and
			\[
				P = \begin{bmatrix} 0 & 1 \\ 1 & 0 \end{bmatrix} \qquad 
				PA = \begin{bmatrix} 2 & 3 \\ 1 & 1 \end{bmatrix}
			\]
			
			\textbf{Factorisation.} Eliminate the $(2,1)$ entry of $PA$ using multiplier $m_{21} = 1/2$:
			\[
				R_2 \leftarrow R_2 - \tfrac12 R_1 = 
				\begin{bmatrix} 1 & 1 \end{bmatrix} - \tfrac12 \begin{bmatrix} 2 & 3 \end{bmatrix} = 
				\begin{bmatrix} 0 & -\tfrac12 \end{bmatrix}
			\]
			This gives
			\[
				L = 
				\begin{bmatrix} 1 & 0 \\ \tfrac12 & 1 \end{bmatrix} \qquad U = 
				\begin{bmatrix} 2 & 3 \\ 0 & -\tfrac12 \end{bmatrix}
			\]
			Check: $LU = \begin{bmatrix} 1 & 0 \\ \tfrac 12 & 1 \end{bmatrix}
			\begin{bmatrix} 2 & 3 \\ 0 & -\tfrac12 \end{bmatrix} = \begin{bmatrix} 2 & 3 \\ 1 & 1 \end{bmatrix} = PA$. Correct.
			
			\textbf{Step 1 (forward substitution):} solve $Ly=Pb$. Since $\mathrm{perm}=[2,1]$, $Pb = 
			\begin{bmatrix} 5 \\ 2 \end{bmatrix}$.
			\[
				y_1 = 5, \qquad y_2 = 2 - \tfrac12(5) = -\tfrac12
			\]
			So $y = \begin{bmatrix}5\\-\tfrac12\end{bmatrix}$.
			
			\textbf{Step 2 (back substitution):} solve $Ux=y$.
			\[
				x_2 = \frac{-\tfrac12}{-\tfrac12} = 1, \qquad x_1 = \frac{1}{2}\left(5 - 3(1)\right) = 1
			\]
			So $x = \begin{bmatrix}1\\1\end{bmatrix}$.
			
			\textbf{Check:} $Ax = 
			\begin{bmatrix} 1 & 1 \\ 2 & 3 \end{bmatrix}
			\begin{bmatrix} 1 \\ 1 \end{bmatrix} = 
			\begin{bmatrix} 2 \\ 5 \end{bmatrix} = b$. Correct.
		\end{example}
		
	\subsection{Partial Pivoting}
		
		\indent \indent The example above shows that the leading-principal-minor condition in Theorem~\ref{thm:lu-existence} can fail in two ways: exactly (Example~\ref{ex:zero-pivot}), or nearly, via a small pivot that leaves the algorithm formally intact but numerically unreliable. \emph{Partial pivoting} is the standard modification to Gaussian elimination that addresses both.
		
		\begin{definition}[Partial Pivoting]
			\label{def:partial-pivoting}
			Let $A \in \mathrm{Mat}(n,\mathbb{F})$. Set $A^{(1)} := A$, and for $k = 2, \dots, n$, let $A^{(k)}$ denote the matrix obtained after the eliminations and row swaps of steps $1,\dots,k-1$ have been applied. At step $k$ of Gaussian elimination, define the \emph{active column} to be the entries $A^{(k)}_{kk}, A^{(k)}_{k+1,k}, \dots, A^{(k)}_{nk}$, i.e.\ column $k$ restricted to rows $k,\dots,n$. Let
			\[
				r := \operatorname*{arg\,max}_{k \le i \le n} 	|A^{(k)}_{ik}|
			\]
			be a row index attaining the largest absolute value in the active column. Partial pivoting swaps row $r$ with row $k$ before proceeding with elimination as normal. Each such swap is a transposition; the composition of these transpositions, in the order performed across steps $k=1,\dots,n-1$, defines the permutation $\sigma \in S_n$ of Definition~\ref{def:permutation-matrix}, and hence the permutation matrix $P$.
		\end{definition}
		
		\noindent With this modification, the substantive question is whether the procedure can always be carried out. That is, whether a valid pivot always exists at every step, for any invertible $A$.
		
		\begin{proposition}[Partial Pivoting on an Invertible Matrix]
			\label{prop:partial-pivoting-exists}
			Let $A \in \mathrm{Mat}(n,\mathbb{F})$ be invertible. Then at every step $k=1,\dots,n-1$ of Gaussian elimination with partial pivoting, the active column contains at least one non-zero entry, so a valid pivot row $r$ always exists.
		\end{proposition}
		\begin{proof}
			A full treatment of this proof is given in \cite{GolubVanLoan2013}.
		\end{proof}
		
		\noindent But why should one bother with partial pivoting? After all, Theorem~\ref{thm:pa-lu-existence} already guarantees some valid $P$ exists for any invertible $A$, so why does the specific pivoting strategy matter?
		
		\begin{proposition}[Bounded Multipliers Under Partial Pivoting]
			\label{prop:bounded-multipliers}
			Suppose partial pivoting is used at every step of Gaussian elimination. Then every multiplier satisfies
			\[
				|m_{ik}| \le 1 \qquad \text{for all } i > k
			\]
		\end{proposition}
		\begin{proof}
			Fix step $k$. By construction, the pivot $A^{(k)}_{kk}$ (after the swap) satisfies
			\[
				|A^{(k)}_{kk}| \ge |A^{(k)}_{ik}| \qquad \text{for 	all } i = k, \dots, n
			\]
			since $A^{(k)}_{kk}$ was chosen to be the entry of largest absolute value among the active column. In particular, for every $i > k$,
			\[
				|m_{ik}| = 	\left|\frac{A^{(k)}_{ik}}{A^{(k)}_{kk}}\right| \le 1
			\]
		\end{proof}
		
	\begin{compremark}
		\texttt{factoriseLU()} (\texttt{lu\_solver.cpp}) makes two representational choices worth justifying against the mathematics of this section.

		\begin{enumerate}[label=\roman*.]
			\item \textbf{In-place storage of the multipliers.} After eliminating column $k$, every entry below the diagonal in that column is exactly zero. Since this value is already known, it carries no information worth keeping, so the position is free to reuse. The code writes the multiplier $m_{ik}$ there instead: $L$'s diagonal is always $1$ (Definition~\ref{def:lu-decomposition}), so nothing about $L$ is lost by leaving it implicit, and the strictly-lower entries of $L$ are exactly these multipliers. So $A$ is overwritten in place to hold both $U$ and the non-trivial part of $L$, with no extra storage needed.
			\item \textbf{The permutation as a vector, not a matrix.} The code applies $P$ to $b$ via \texttt{bPrime(i) = b(perm[i], 0)}, a direct reindexing. This matches how $P$ is actually used in the proof of Proposition~\ref{prop:solve-pa-lu}: the argument never needs more from $P$ than reordering $b$'s rows. So representing $P$ as the vector \texttt{perm}, rather than an $n\times n$ matrix, is a faithful, more efficient match for what the proof uses.
		\end{enumerate}
		
		Both choices can be checked directly against the worked example (Example~\ref{ex:solve-pa-lu-nontrivial}): running \texttt{lu\_solver.cpp} on that example's $A$ and $b$ reproduces $x = (1,1)^\top$.
	\end{compremark}
	
\section{Inner Product Spaces and QR Decompositions}
\label{sec:inner-products}
	
	This section develops inner product spaces, orthogonality and the QR decomposition, providing the tools to construct orthonormal bases and to solve $Ax = b$ via orthogonal factorisation. \\
	
	Beyond providing a second, more numerically stable route to solving linear systems, orthogonality is the structural idea that the spectral theorem in the next section depends on directly. \\
	
	The algorithm is implemented from scratch in C++ as \texttt{qr\_gram\_schmidt.cpp} (Appendix~\ref{app:qr-solver}), computing the QR decomposition via the Gram-Schmidt process.
	
	\subsection{Inner Product Spaces}
		
		\indent \indent In Section~\ref{sec:vector-spaces-and-linear-operators}, we defined a vector space $V$, but a vector space alone carries no notion of length or angle between vectors. Equipping $V$ with such a notion is the first step toward orthogonality, which in turn is the basis of the QR decomposition developed later in this section.
		
		Recall the dot product for $\mathbb{R}^n$:
		\[
			x \cdot y = \sum_{i=1}^{n} x_i y_i
		\]
		
		This immediately gives a notion of length via $\|x\| := \sqrt{x \cdot x}$.
		
		The dot product on $\mathbb{R}^n$ has three properties worth isolating: it is symmetric in $x$ and $y$, linear in each argument and positive-definite ($x \cdot x \ge 0$, with equality iff $x=0$). Abstracting exactly these properties to an arbitrary vector space yields the definition of an inner product.
		
		Throughout this section, $V$ will denote a real vector space, i.e.\ $\mathbb{F} = \mathbb{R}$.
		
		\begin{definition}[Inner Product / Inner Product Space]
			\label{def:inner-product}
			An inner product on a real vector space $V$ is a map $\langle \cdot, \cdot \rangle : V \times V \to \mathbb{R}$ that satisfies the following properties for all $u, v, w \in V$ and all $\lambda \in \mathbb{R}$:
			
			\begin{enumerate}[label=(IP\arabic*)]
				\item \label{ax:IP1} \textbf{Symmetry:} $\langle u, v \rangle = \langle v, u \rangle$
				\item \label{ax:IP2} \textbf{Linearity in the first argument:} $\langle u + v, w \rangle = \langle u, w \rangle + \langle v, w \rangle$ and $\langle \lambda u, v \rangle = \lambda \langle u, v \rangle$
				\item \label{ax:IP3} \textbf{Positive-definiteness:} $\langle v, v \rangle \ge 0$ for all $v \in V$, with equality $\langle v, v \rangle = 0$ if and only if $v = 0_V$
			\end{enumerate}
			
			The pair $(V, \langle \cdot, \cdot \rangle)$ is called an \textbf{inner product space}.
			
		\end{definition}
		
		\begin{example}[Trace Inner Product]
			\label{ex:trace-inner-product}
			Consider $\mathrm{Mat}(m, n, \mathbb{R})$. Equipped with the trace inner product $\langle A | B \rangle := \operatorname{tr}(A^\top B)$ for all $A, B \in \mathrm{Mat}(m, n, \mathbb{R})$ the pair $(\mathrm{Mat}(m, n, \mathbb{R}), \langle \cdot, \cdot \rangle)$ is an inner product space. \\
			\indent Further, by Definition~\ref{def:induced-norm}, this inner product's induced norm is $\|A\|_F := \sqrt{\langle A,A\rangle} = \sqrt{\operatorname{tr}(A^\top A)}$, called the Frobenius norm.
		\end{example}
		
		\begin{remark}
			In a real vector space, the properties of symmetry \ref{ax:IP1} and linearity in the first argument \ref{ax:IP2} together imply linearity in the second argument, meaning the inner product is fully bilinear.
		\end{remark} 
		
		\begin{definition}[Induced Norm] 
			\label{def:induced-norm}
			Let $(V, \langle \cdot, \cdot \rangle)$ be an inner product space. For each vector $v \in V$, the \textbf{norm} of $v$ induced by the inner product, denoted by $\|v\|$, is defined as:
			\[
				\|v\| := \sqrt{\langle v, v \rangle} 
			\]
		\end{definition}
		
		\begin{remark}
			\label{rem:bilinearity}
			Note that the quantity ${\langle v, v \rangle}$ is non-negative from positive-definiteness \ref{ax:IP3}, hence the induced norm is always well-defined in $\mathbb{R}$.
		\end{remark}
		
		\begin{theorem}[Cauchy-Schwarz Inequality] \label{thm:cauchy-schwarz}
			Let $V$ be an inner product space. For all vectors $u, v \in V$, the following inequality holds:
			\[ 
				|\langle u, v \rangle| \le \|u\| \|v\|
			\]
			This inequality is an equality if and only if $u, v$ are linearly dependent (Definition~\ref{def:lin-dep}).
		\end{theorem}
		
		\begin{proof}
			Let $V$ be a real inner product space. We wish to show that for all $u, v \in V$, $|\langle u, v \rangle| \le \|u\| \|v\|$.
			
			Suppose $v = 0_V$. $0_V = 0 \cdot u$ (Proposition~\ref{prop:vspace-properties}); combined with symmetry \ref{ax:IP1} and linearity in the first argument \ref{ax:IP2}, $\langle u, 0_V \rangle = \langle u, 0 \cdot u \rangle = \langle 0\cdot u, u\rangle = 0\langle u,u\rangle = 0$. The inequality then becomes:
			\[ 
				|\langle u, 0_V \rangle| = 0 \le \|u\| \cdot 0 = 0
			\]
			which holds as an equality. Thus, the theorem is true for $v = 0_V$.
			
			Now assume $v \neq 0_V$. For any real number $t \in \mathbb{R}$, consider the function:
			\[ 
				f(t) = \|u - tv\|^2
			\]
			By the definition of the induced norm, $f(t) = \langle u - tv, u - tv \rangle$. By positive-definiteness \ref{ax:IP3}, the inner product of any vector with itself is non-negative, and is zero iff the vector is the zero vector. Therefore, $f(t) \ge 0$ for all $t \in \mathbb{R}$.
			
			Using bilinearity (Remark~\ref{rem:bilinearity}), we expand the inner product:
			
			\begin{align*}
				f(t) &= \langle u, u - tv \rangle - t\langle v, u - tv \rangle \\
				&= \langle u, u \rangle - t\langle u, v \rangle - t\langle v, u \rangle + t^2\langle v, v \rangle \\
				&= \|u\|^2 - 2t\langle u, v \rangle + t^2\|v\|^2
			\end{align*}
			
			This is a quadratic in $t$ of the form $at^2 + bt + c$, where $a = \|v\|^2$, $b = -2\langle u, v \rangle$, and $c = \|u\|^2$.
			
			Since $v \neq 0_V$, positive-definiteness \ref{ax:IP3} gives $\|v\|^2 = \langle v,v\rangle > 0$, so $\|v\|^2 \neq 0$ and division by it below is valid.
			
			Following the minimisation route, we choose the specific value of $t$ that minimises $f(t)$. For a quadratic $at^2 + bt + c$, the minimum occurs at $t = -b/(2a)$. Here, we set:
			\[
				t_0 = \frac{2\langle u, v \rangle}{2\|v\|^2} = \frac{\langle u, v \rangle}{\|v\|^2}
			\]
			Plugging $t_0$ into our inequality $f(t) \ge 0$:
			
			\begin{align*}
				0 &\le \|u\|^2 - 2\left(\frac{\langle u, v \rangle}{\|v\|^2}\right)\langle u, v \rangle + \left(\frac{\langle u, v \rangle}{\|v\|^2}\right)^2 \|v\|^2 \\
				0 &\le \|u\|^2 - \frac{2\langle u, v \rangle^2}{\|v\|^2} + \frac{\langle u, v \rangle^2}{\|v\|^2} \\
				0 &\le \|u\|^2 - \frac{\langle u, v \rangle^2}{\|v\|^2}
			\end{align*}
			
			Rearranging terms, we find $\frac{\langle u, v \rangle^2}{\|v\|^2} \le \|u\|^2$, which implies $\langle u, v \rangle^2 \le \|u\|^2 \|v\|^2$. Taking the square root of both sides yields the desired inequality:
			\[ 
				|\langle u, v \rangle| \le \|u\| \|v\|
			\]
			
			Equality $|\langle u, v \rangle| = \|u\| \|v\|$ holds iff $f(t_0) = 0$. By positive-definiteness \ref{ax:IP3}, $f(t_0) = \|u - t_0v\|^2 = 0$ iff $u - t_0v = 0_V$. This implies $u = t_0v$, meaning $u$ is a scalar multiple of $v$. Thus, equality holds iff $u$ and $v$ are linearly dependent.
		\end{proof}
		
		\begin{remark}[Angle Between Vectors]
			\label{rem:angle-between-vectors}
			Let $V$ be a real inner product space and let $u, v \in V$ with $u, v \neq 0_V$. By Cauchy-Schwarz (Theorem~\ref{thm:cauchy-schwarz}),
			\[
				|\langle u, v \rangle| \le \|u\|\|v\|
			\]
			Since $u, v \neq 0_V$, positive-definiteness \ref{ax:IP3} gives $\|u\|^2 > 0$ and $\|v\|^2 > 0$, so $\|u\| \neq 0$ and $\|v\| \neq 0$, and we may divide through by $\|u\|\|v\|$:
			\[
				\left| \frac{\langle u,v \rangle}{\|u\|\|v\|} \right| \le 1
			\]
			That is, the ratio $\frac{\langle u, v \rangle}{\|u\|\|v\|}$ always lies in $[-1,1]$, which is exactly the range of $\cos$. Since $\cos : [0,\pi] \to [-1,1]$ is a bijection, there is a unique $\theta \in [0,\pi]$ satisfying
			\[
				\cos\theta = \frac{\langle u,v \rangle}{\|u\|\|v\|}
			\]
			This uniqueness is precisely what licenses calling $\theta$ \textbf{the angle between} $u$ and $v$: without Cauchy-Schwarz bounding the ratio into $[-1,1]$ first, no such $\theta$ would be well-defined at all.
		\end{remark}
		
		\begin{proposition}[Triangle Inequality]
			\label{prop:triangle-inequality}
			Let $V$ be a real inner product space. Then for all $u, v \in V$
			\[
				\|u+v\| \le \|u\| + \|v\|
			\]
			(a consequence of Cauchy-Schwarz, Theorem~\ref{thm:cauchy-schwarz}).
		\end{proposition}
		
		\begin{proof}
			We bound $\|u+v\|^2$ and take square roots at the end.
			
			By the definition of the induced norm and bilinearity (Remark~\ref{rem:bilinearity}),
			\begin{align*}
				\|u+v\|^2 &= \langle u+v, u+v \rangle \\
				&= \langle u,u\rangle + \langle u,v\rangle + \langle v,u\rangle + \langle v,v\rangle \\
				&= \|u\|^2 + 2\langle u,v\rangle + \|v\|^2
			\end{align*}
			using symmetry \ref{ax:IP1} to combine $\langle u,v\rangle$ and $\langle v,u\rangle$.
			
			By Cauchy-Schwarz (Theorem~\ref{thm:cauchy-schwarz}), $|\langle u,v\rangle| \le \|u\|\|v\|$, and in particular $\langle u,v\rangle \le |\langle u,v\rangle| \le \|u\|\|v\|$. Substituting this bound on the cross term,
			\[
				\|u+v\|^2 \le \|u\|^2 + 2\|u\|\|v\| + \|v\|^2
			\]
			The right-hand side is a perfect square:
			\[
				\|u\|^2 + 2\|u\|\|v\| + \|v\|^2 = \left(\|u\| + \|v\|\right)^2
			\]
			So $\|u+v\|^2 \le \left(\|u\|+\|v\|\right)^2$. Since both $\|u+v\|$ and $\|u\|+\|v\|$ are non-negative, taking square roots preserves the inequality:
			\[
				\|u+v\| \le \|u\| + \|v\|
			\]
		\end{proof}
		
	\subsection{Orthogonality}
		
		\noindent The angle between vectors (Remark~\ref{rem:angle-between-vectors}) showed that $\theta = \pi/2$ if and only if $\cos\theta = 0$, which by that same remark holds if and only if $\langle u, v \rangle = 0$. This special case (i.e. an inner product of zero) is important enough to warrant its own name.
		
		\begin{definition}[Orthogonal Set / Orthonormal Set]
			\label{def:orthogonal-orthonormal-set}
			Let $V$ be an inner product space and let $v_1, \dots, v_m \in V$ be a set of vectors.
			\begin{enumerate}[label=\roman*.]
				\item The set $\{v_1,\dots,v_m\}$ is \emph{orthogonal} if $\langle v_i, v_j \rangle = 0$ for all $i \neq j$.
				\item The set $\{v_1,\dots,v_m\}$ is \emph{orthonormal} if it is orthogonal and, in addition, $\|v_i\| = 1$ for all $i$.
			\end{enumerate}
		\end{definition}
		
		\begin{definition}[Orthogonal Matrix]
			\label{def:orthogonal-matrix}
			A matrix $Q \in \mathrm{Mat}(n,\mathbb{R})$ is \emph{orthogonal} if its columns form an orthonormal set (Definition~\ref{def:orthogonal-orthonormal-set}), equivalently, if
			\[
			Q^\top Q = \mathbbm{1}
			\]
		\end{definition}
		
		\begin{remark}
			Both conditions in Definition~\ref{def:orthogonal-orthonormal-set} compress into a single statement using the Kronecker delta: $\{v_1,\dots,v_m\}$ is orthonormal if and only if
			\[
				\langle v_i, v_j \rangle = \delta_{ij} \qquad \text{for all } i,j
			\]
			where $\delta_{ij} = 1$ if $i=j$ and $0$ otherwise.
		\end{remark}
		
		\begin{proposition}[Orthogonality and Linear Independence]
			\label{prop:orthogonal-and-linearly-independent}
			Let $V$ be an inner product space and let $v_1, \dots, v_m \in V$ be an orthogonal set of nonzero vectors. Then $v_1, \dots, v_m$ are linearly independent.
		\end{proposition}
		
		\begin{proof}
			Suppose
			\[
				c_1 v_1 + \cdots + c_m v_m = 0_V
			\]
			for some scalars $c_1,\dots,c_m \in \mathbb{R}$. We show $c_i = 0$ for all $i$.
			
			Fix $k \in \{1,\dots,m\}$ and take the inner product of both sides with $v_k$:
			\[
				\left\langle \sum_{i=1}^{m} c_i v_i,\; v_k \right\rangle = \langle 0_V, v_k \rangle
			\]
			By symmetry \ref{ax:IP1} and the same argument as in the Cauchy–Schwarz proof, $\langle 0_V, v_k \rangle = \langle v_k, 0_V \rangle = 0$.
			
			Expanding the left-hand side by bilinearity (Remark~\ref{rem:bilinearity}),
			\[
				\sum_{i=1}^{m} c_i \langle v_i, v_k \rangle = 0
			\]
			Since $\{v_1,\dots,v_m\}$ is orthogonal, $\langle v_i, v_k\rangle = 0$ for all $i \neq k$, so every term with $i\neq k$ vanishes, leaving only the $i=k$ term:
			\[
				c_k \langle v_k, v_k \rangle = 0, \qquad \text{i.e.} \qquad c_k \|v_k\|^2 = 0
			\]
			Since $v_k \neq 0_V$, positive-definiteness \ref{ax:IP3} gives $\|v_k\|^2 > 0$, so $\|v_k\|^2 \neq 0$, and hence $c_k = 0$.
			
			Since $k$ was arbitrary, $c_k = 0$ for every $k=1,\dots,m$, so $v_1,\dots,v_m$ are linearly independent.
		\end{proof}
		
		\begin{theorem}[Gram-Schmidt Process]
			\label{thm:gram-schmidt}
			Let $V$ be an inner product space and let $v_1, \dots, v_m \in V$ be linearly independent. Define
			\[
				u_1 = v_1, \qquad u_k = v_k - \sum_{j=1}^{k-1} \frac{\langle v_k, u_j \rangle}{\langle u_j, u_j \rangle} u_j \quad \text{for } k = 2, \dots, m
			\]
			Then $u_1, \dots, u_m$ is an orthogonal, hence linearly independent (Proposition~\ref{prop:orthogonal-and-linearly-independent}), set of nonzero vectors satisfying
			\[
				\operatorname{span}\{u_1, \dots, u_k\} = \operatorname{span}\{v_1, \dots, v_k\} \qquad \text{for every } k = 1, \dots, m
			\]
		\end{theorem}
		
		\begin{proof}
			A full treatment of this proof is given in \cite{GolubVanLoan2013}.
		\end{proof}
		
		\begin{corollary}[Existence of an Orthonormal Basis]
			\label{cor:orthonormal-basis-exists}
			Let $V$ be a finite-dimensional inner product space with $\dim V = n$. Then, by Theorem~\ref{thm:gram-schmidt}, $V$ has an orthonormal basis.
		\end{corollary}
		
		\begin{proof}
			Since $\dim V = n$, $V$ has a basis $v_1,\dots,v_n$.
			
			Apply Gram-Schmidt (Theorem~\ref{thm:gram-schmidt}) to $v_1,\dots,v_n$, producing an orthogonal set $u_1,\dots,u_n$ with
			\[
				\operatorname{span}_{\mathbb{R}}\{u_1,\dots,u_n\} = \operatorname{span}_{\mathbb{R}}\{v_1,\dots,v_n\} = V
			\]
			By Theorem~\ref{thm:gram-schmidt}, $u_1,\dots,u_n$ is an orthogonal, linearly independent set of nonzero vectors, with nonzero vectors. Being $n$ linearly independent vectors spanning the $n$-dimensional space $V$, $u_1,\dots,u_n$ form a basis of $V$ by Proposition~\ref{prop:v-basis}.
			
			Define $e_k := u_k / \|u_k\|$ for each $k=1,\dots,n$. Since $u_k \neq 0_V$, positive-definiteness \ref{ax:IP3} gives $\|u_k\| \neq 0$, so this division is well-defined.
			
			For $i \neq k$, by bilinearity (Remark~\ref{rem:bilinearity}),
			\[
				\langle e_i, e_k \rangle = \left\langle \frac{u_i}{\|u_i\|}, \frac{u_k}{\|u_k\|} \right\rangle = \frac{1}{\|u_i\|\|u_k\|}\langle u_i, u_k \rangle = \frac{1}{\|u_i\|\|u_k\|} \cdot 0 = 0
			\]
			since $\langle u_i, u_k\rangle = 0$ by orthogonality of $u_1,\dots,u_n$, and dividing zero by the nonzero scalar $\|u_i\|\|u_k\|$ leaves zero.
			
			By bilinearity (Remark~\ref{rem:bilinearity}),
			\[
				\|\lambda v\|^2 = \langle \lambda v, \lambda v \rangle = \lambda^2 \langle v,v \rangle = \lambda^2 \|v\|^2
			\]
			so $\|\lambda v\| = |\lambda| \, \|v\|$ (taking non-negative square roots). Applying this with $\lambda = 1/\|u_k\|$,
			\[
				\|e_k\| = \left\| \frac{u_k}{\|u_k\|} \right\| = \frac{1}{\|u_k\|}\|u_k\| = 1
			\]
			
			Hence $\{e_1,\dots,e_n\}$ is orthonormal (Definition~\ref{def:orthogonal-orthonormal-set}). Since scaling each $u_k$ by the nonzero scalar $1/\|u_k\|$ changes neither the span nor linear independence of $\{u_1,\dots,u_n\}$, $\{e_1,\dots,e_n\}$ remains a basis of $V$. Thus $\{e_1,\dots,e_n\}$ is an orthonormal basis of $V$.
		\end{proof}
		
		\begin{remark}[Gram-Schmidt Fixes an Already-Orthonormal Set]
			\label{rem:gram-schmidt-fixes-orthonormal}
			If $u_1,\dots,u_r$ are already orthonormal, feeding them as the first $r$ inputs to Gram-Schmidt (Theorem~\ref{thm:gram-schmidt}) changes nothing: at step $k \le r$, the formula
			\[
				u_k - \sum_{j=1}^{k-1} \frac{\langle u_k, u_j \rangle}{\langle u_j, u_j \rangle} u_j
			\]
			subtracts off projections onto $u_1,\dots,u_{k-1}$, but these are all zero since the $u_j$ are already mutually orthogonal — so the sum vanishes and $u_k$ passes through unchanged. Since each $u_k$ is already a unit vector, the normalisation $e_k := u_k/\|u_k\|$ (as in Corollary~\ref{cor:orthonormal-basis-exists}) leaves it unchanged too. So Gram-Schmidt only starts doing real work from step $r+1$ onward.
		\end{remark}
		
	\subsection{QR Decompositions}
		
		\begin{definition}[QR Decomposition]
			\label{def:qr-decomp}
			Let $A \in \mathrm{Mat}(n, \mathbb{R})$. A QR decomposition of $A$ is a factorisation $A=QR$ where $Q \in \mathrm{Mat}(n,\mathbb{R})$ has orthonormal columns, that is, the columns $q_1, \dots, q_n$ form an orthonormal set (Definition~\ref{def:orthogonal-orthonormal-set}); and $R \in \mathrm{Mat}(n,\mathbb{R})$ is upper triangular with positive diagonal entries.
		\end{definition}
		
		\begin{theorem}[Existence and Uniqueness of the QR Decomposition]
			\label{thm:existence-qr-decomp}
			Let $A \in \mathrm{Mat}(n,\mathbb{R})$ be invertible. Then $A$ admits a QR decomposition $A=QR$ (as in Definition~\ref{def:qr-decomp}), and this decomposition is unique.
		\end{theorem}
		
		\begin{proof}
			Let $a_1,\dots, a_n$ denote the columns of $A$. Since $A$ is invertible, its columns are linearly independent.
			
			Apply Gram-Schmidt (Theorem~\ref{thm:gram-schmidt}) to $a_1,\dots,a_n$ to obtain an orthogonal set $u_1,\dots,u_n$ with
			\[
				\operatorname{span}_{\mathbb{R}}\{u_1,\dots,u_k\} = \operatorname{span}_{\mathbb{R}}\{a_1,\dots,a_k\} \qquad \text{for every } k=1, \dots, n
			\]
			and normalise $q_k := u_k/\|u_k\|$ for each $k$, exactly as in the construction of Corollary~\ref{cor:orthonormal-basis-exists}; we do not re-derive well-definedness or orthonormality of $q_1,\dots,q_n$ here, having already established both.
			
			Let $Q$ be the matrix with columns $q_1,\dots,q_n$. By construction, $Q$ has orthonormal columns.
			
			We now construct $R$. Since $\operatorname{span}_{\mathbb{R}}\{q_1,\dots,q_k\} = \operatorname{span}_{\mathbb{R}}\{u_1,\dots,u_k\} = \operatorname{span}_{\mathbb{R}}\{a_1,\dots,a_k\}$ for every $k$, each $a_k$ lies in $\operatorname{span}_{\mathbb{R}}\{q_1,\dots,q_k\}$, so
			\[
			a_k = \sum_{j=1}^{k} c_j\, q_j
			\]
			for some scalars $c_1,\dots,c_k$. Fix $i \le k$ and take the inner product of both sides with $q_i$. By bilinearity (Remark~\ref{rem:bilinearity}),
			\[
				\langle a_k, q_i \rangle = \sum_{j=1}^{k} c_j \langle q_j, q_i \rangle
			\]
			Since $\{q_1,\dots,q_n\}$ is orthonormal, $\langle q_j,q_i\rangle = 0$ for $j\neq i$ and $\langle q_i,q_i\rangle = 1$, so every term vanishes except $j=i$, leaving $\langle a_k, q_i\rangle = c_i$. As $i \le k$ was arbitrary, this identifies every coefficient: $c_j = \langle a_k, q_j\rangle$ for all $j=1,\dots,k$, so
			\[
				a_k = \sum_{j=1}^{k} \langle a_k, q_j \rangle\, q_j
			\]
			Define $R_{jk} := \langle a_k, q_j \rangle$ for $j \le k$ and $R_{jk} := 0$ for $j > k$. This is upper triangular by construction, and the identity above shows precisely that $a_k = \sum_{j=1}^n R_{jk} q_j$ is the $k$-th column of $QR$, so $A = QR$.
			
			It remains to show $R_{kk} > 0$ for every $k$. Fix $k$. By the Gram-Schmidt relation,
			\[
				u_k = a_k - \sum_{j<k} \operatorname{proj}_{u_j}(a_k)
			\]
			so
			\[
				a_k = u_k + \sum_{j<k} \operatorname{proj}_{u_j}(a_k)
			\]
			Each $\operatorname{proj}_{u_j}(a_k)$ is by definition a scalar multiple of $u_j$, hence of $q_j$. So $a_k = u_k + \sum_{j<k} d_j q_j$ for some scalars $d_j$. Taking the inner product with $q_k$ and using bilinearity (Remark~\ref{rem:bilinearity}),
			\[
				\langle a_k, q_k \rangle = \langle u_k, q_k \rangle + \sum_{j<k} d_j \langle q_j, q_k \rangle
			\]
			Since $\{q_1,\dots,q_n\}$ is orthonormal, $\langle q_j,q_k\rangle = 0$ for $j \neq k$, so every term in the sum vanishes, leaving
			\[
				R_{kk} = \langle a_k, q_k \rangle = \langle u_k, q_k \rangle = \left\langle u_k, \frac{u_k}{\|u_k\|} \right\rangle = \frac{\|u_k\|^2}{\|u_k\|} = \|u_k\|
			\]
			Since $u_k \neq 0_V$ (Gram-Schmidt produces nonzero vectors, as used already in Corollary~\ref{cor:orthonormal-basis-exists}), positive-definiteness \ref{ax:IP3} gives $\|u_k\| > 0$. Hence $R_{kk} > 0$ for all $k$. This establishes existence of the decomposition $A = QR$ with the stated properties.
			
			To prove uniqueness of the QR decomposition, we rely on \cite{GolubVanLoan2013}.
		\end{proof}
		
		\begin{proposition}[Solving $Ax=b$ via $A=QR$]
			\label{prop:solve-qr}
			Let $A \in \mathrm{Mat}(n,\mathbb{R})$ be invertible, let $Q, R$ be such that $A = QR$ (Theorem~\ref{thm:existence-qr-decomp}), and let $b \in \mathbb{R}^n$. Then the system $Ax=b$ has a unique solution $x$, obtained by first computing $c := Q^\top b$, then solving $Rx = c$ for $x$ by back substitution.
		\end{proposition}
		
		\begin{proof}
			Substituting $A = QR$, the system $Ax=b$ becomes
			\[
				QRx = b
			\]
			Since $Q$ has orthonormal columns, $Q^\top Q = I$. As $Q$ is square, this left inverse is automatically a two-sided inverse, so $Q$ is invertible with $Q^{-1} = Q^\top$.
			
			Multiplying both sides of $QRx=b$ on the left by $Q^\top$,
			\[
				Q^\top Q R x = Q^\top b \quad \Longrightarrow \quad Rx = Q^\top b
			\]
			Define $c := Q^\top b$. Unlike the analogous step in the $PA=LU$ method, computing $c$ requires no triangular solve — it is a single direct matrix-vector product.
			
			It remains to solve $Rx = c$. Since $R$ is upper triangular with $R_{kk} > 0$ for all $k$ (Theorem~\ref{thm:existence-qr-decomp}), $R$ is invertible, so $Rx=c$ has a unique solution obtained by back substitution:
			\[
				x_k = \frac{1}{R_{kk}}\left(c_k - \sum_{j=k+1}^{n} R_{kj}x_j\right) \qquad k = n, n-1, \dots, 1
			\]
			Existence and uniqueness of $c$ and then $x$ give existence and uniqueness of a solution to $Ax=b$.
		\end{proof}
		
	\begin{compremark}
		\texttt{gramSchmidtQR()} (Appendix~\ref{app:qr-solver}) projects each $R_{jk}$ against the running vector $u_k$ and not the original column $a_k$ as Theorem~\ref{thm:gram-schmidt}'s formula suggests. This is fine mathematically: once earlier projections are removed from $u_k$, $\langle u_k,q_j\rangle = \langle a_k,q_j\rangle$ anyway, so the theorem's guarantees still hold. This is the modified Gram-Schmidt Process.
		
		It matters numerically because projecting against $a_k$ every time (classical Gram-Schmidt) lets rounding error compound independently in each coefficient; projecting against the already-partially-cleaned $u_k$ (modified Gram-Schmidt) stops that error from carrying forward as badly. Refer to \cite[\S 5.2.7--5.2.9]{GolubVanLoan2013} for a more detailed explanation.
		
		This is verified directly in \texttt{main()}: $Q^\top Q \approx \mathbbm{1}$ to machine precision ($10^{-16}$--$10^{-17}$ off the diagonal), and $QR \approx A$, an exact reconstruction of the input.
	\end{compremark}
	
\section{Eigenvalues and the Spectral Theorem}
\label{sec:eigenvalues-spectral-theorem}
	
	This section develops eigenvalues, eigenvectors and the spectral theorem for symmetric matrices, culminating in the existence of an orthonormal eigenbasis. \\
	
	Beyond its role in diagonalising symmetric matrices, this result is the algebraic foundation for the Singular Value Decomposition (or SVD), which this paper uses later in the chapter to analyse and decompose general $m\times n$ matrices. \\
	
	The algorithm is implemented from scratch in C++ as \texttt{jacobi\_eigensolver.cpp} (Appendix~\ref{app:jacobi-solver}), computing the eigendecomposition of a symmetric matrix via the classical Jacobi method.
	
	\subsection{Eigenvalues and Eigenvectors}
	\label{sec:eigenvalues-and-eigenvectors}
		
		Let $A \in \mathrm{Mat}(n,\mathbb{R})$.
		
		\begin{definition}[Eigenvalue / Eigenvector]
			\label{def:eigenvalue-eigenvector}
			A scalar $\lambda \in \mathbb{F}$ is an eigenvalue of $A$ if there exists a nonzero $v \in \mathbb{F}^n$ with $Av = \lambda v$. Such a $v$ is an eigenvector.
		\end{definition}
		
		Rewriting the above statement's equation, we obtain $Av = \lambda v \iff (A - \lambda \mathbbm{1})v = 0$. So $\lambda$ is an eigenvalue iff $A - \lambda\mathbbm{1}$ is singular, i.e.
		\[
			\det(A - \lambda\mathbbm{1}) = 0
		\]
		
		\noindent This constitutes the following definition.
		
		\begin{definition}[Characteristic Polynomial]
			\label{def:characteristic-polynomial}
			The characteristic polynomial of a matrix $A$ is a degree-$n$ polynomial in $\lambda$ of the form
			\[
				p_A(\lambda) := \det(A - \lambda\mathbbm{1})
			\]
			Its roots are exactly the eigenvalues of $A$.
		\end{definition}
		
		\begin{remark}
			\label{rem:existence-eigenvalues-complex}
			Over $\mathbb{C}$, the Fundamental Theorem of Algebra guarantees that $p_A$ has exactly $n$ roots counted with multiplicity, so every $A \in \mathrm{Mat}(n,\mathbb{C})$ has $n$ eigenvalues.
		\end{remark}
		
		\begin{example}[A Matrix with No Real Eigenvalues]
			\label{ex:no-real-eigenvalues}
			Consider
			\[
				A = \begin{bmatrix} 0 & -1 \\ 1 & 0 \end{bmatrix}
			\]
			The characteristic polynomial is
			\[
				p_A(\lambda) = \det(A - \lambda \mathbbm{1}) = \det \begin{bmatrix} -\lambda & -1 \\ 1 & -\lambda \end{bmatrix} = (-\lambda)(-\lambda) - (-1)(1) = \lambda^2 + 1
			\]
			Setting $p_A(\lambda) = 0$ gives $\lambda^2 = -1$, which has no solution in $\mathbb{R}$; the roots are $\lambda = \pm i$. So $A$ has no real eigenvalues, and hence no real eigenvectors.
			
			Geometrically, $A$ is the matrix of a $90°$ rotation (anticlockwise) in $\mathbb{R}^2$. An eigenvector would be a nonzero vector whose direction is unchanged by $A$ but a $90°$ rotation fixes no direction at all: every nonzero vector in $\mathbb{R}^2$ is rotated onto a different line through the origin. This is the geometric content behind the absence of real eigenvalues here.
		\end{example}
		
		This motivates restricting our attention to a class of matrices for which real eigenvalues are guaranteed: symmetric matrices.
		
		\begin{definition}[Symmetric Matrix]
			\label{def:symmetric-matrix}
			Let $A \in \mathrm{Mat}(n, \mathbb{R})$. A symmetric matrix is a matrix where $A^\top = A$.
		\end{definition}
		
		\begin{definition}[Self Adjoint]
			\label{def:self-adjoint}
			Let $A \in \mathrm{Mat}(n,\mathbb{R})$, equipped with the standard inner product. When the following property is true for $A$,
			\[
				\langle Av,w\rangle = \langle v,Aw\rangle \quad \text{for all } v,w \in \mathbb{R}^n
			\]
			$A$ is called self-adjoint.
		\end{definition}
		
		\begin{remark}
			\label{rem:self-adjoint}
			Some remarks to consider about self-adjointness are
			\begin{enumerate}[label=\roman*.]
				\item When $A$ is symmetric, the property holds.
				\item The property in Definition~\ref{def:self-adjoint} can be generalised for any inner product, however for the scope of this paper we define this one case.
			\end{enumerate}
		\end{remark}
		
		\begin{theorem}[Real Eigenvalues of Symmetric Matrices]
			\label{thm:symmetric-real-eigenvalues}
			Let $A \in \mathrm{Mat}(n,\mathbb{R})$ be symmetric, i.e.\ $A = A^\top$. Then every eigenvalue of $A$ is real.
		\end{theorem}
		
		\begin{proof}
			Let $\lambda \in \mathbb{C}$ be a root of the characteristic polynomial $p_A$, with corresponding eigenvector $v \in \mathbb{C}^n \setminus \{0\}$, so $Av = \lambda v$; a priori $\lambda$ and $v$ may be complex, since we have not yet shown $\lambda$ is real. Let $v^* := \bar{v}^\top$ denote the conjugate transpose of $v$, and consider the scalar $v^*Av$.
			
			Using $Av = \lambda v$,
			\[
				v^*Av = v^*(\lambda v) = \lambda\, v^*v
			\]
			
			We show $v^*Av$ is real, by showing it equals its own complex conjugate. Conjugating the product $v^*Av = \bar v^\top A v$ entrywise (conjugation of a product preserves the order of the factors, unlike transpose),
			\[
				\overline{v^*Av} = \overline{\bar v^\top A v} = \overline{\bar v}^{\,\top} \bar A\, \bar v = v^\top A \bar v
			\]
			using $\overline{\bar v} = v$ and, since $A$ is real, $\bar A = A$.
			
			Separately, since $v^*Av$ is a $1\times 1$ matrix (a scalar), it equals its own transpose, and transposing a product reverses the order of the factors:
			\[
				v^*Av = \bar v^\top A v = \left(\bar v^\top A v\right)^\top = v^\top A^\top \bar v = v^\top A \bar v
			\]
			using symmetry, $A^\top = A$, in the last step.
			
			Both computations give $v^\top A \bar v$, so $\overline{v^*Av} = v^*Av$, and hence $v^*Av$ is real.
			
			Writing $v = (v_1,\dots,v_n)^\top$,
			\[
				v^*v = \bar v^\top v = \sum_{i=1}^{n} \bar v_i v_i = \sum_{i=1}^{n} |v_i|^2
			\]
			which is real, and strictly positive since $v \neq 0$.
			
			From the first computation, $v^*Av = \lambda\, v^*v$. The left-hand side is real, and $v^*v$ is a positive real number; a real number divided by a positive real number is real, so
			\[
				\lambda = \frac{v^*Av}{v^*v}
			\]
			is real.
		\end{proof}
		
		\begin{remark}[Standard Inner Product]
			\label{rem:inner-product-as-matmul}
			On $\mathbb{R}^n$, the standard inner product coincides with the dot product, $\langle x,y \rangle := x \cdot y = \sum_{i=1}^n x_iy_i$, which for column vectors $x,y \in \mathbb{R}^n$ can equivalently be written as the matrix product
			\[
				\langle x,y \rangle = x^\top y
			\]
			a $1\times 1$ matrix identified with its unique scalar entry. This identification is used throughout the results on symmetric matrices below.
		\end{remark}
		
		\begin{proposition}[Orthogonal Eigenvectors]
			\label{prop:symmetric-eigenvectors-orthogonal}
			Let $A \in \mathrm{Mat}(n,\mathbb{R})$ be symmetric, and let $v_1, v_2$ be eigenvectors of $A$ with $Av_1 = \lambda_1 v_1$, $Av_2 = \lambda_2 v_2$, and $\lambda_1 \neq \lambda_2$. Then $\langle v_1, v_2 \rangle = 0$.
		\end{proposition}
		
		\begin{proof}
			Consider $\langle Av_1, v_2 \rangle$, computed two ways.
			
			Using $Av_1 = \lambda_1 v_1$ and linearity in the first argument (Remark~\ref{rem:bilinearity}),
			\[
				\langle Av_1, v_2 \rangle = \langle \lambda_1 v_1, v_2 \rangle = \lambda_1 \langle v_1, v_2 \rangle
			\]
			
			Using the matrix-multiplication form of the inner product (Remark~\ref{rem:inner-product-as-matmul}), $\langle Av_1,v_2\rangle = (Av_1)^\top v_2 = v_1^\top A^\top v_2$, and symmetry, $A^\top = A$,
			\[
				\langle Av_1, v_2 \rangle = v_1^\top A^\top v_2 = v_1^\top A v_2 = \langle v_1, Av_2 \rangle = \langle v_1, \lambda_2 v_2 \rangle = \lambda_2 \langle v_1, v_2 \rangle
			\]
			where the last step uses linearity in the second argument (Remark~\ref{rem:bilinearity}).
			
			Equating the two computations,
			\[
				\lambda_1 \langle v_1, v_2 \rangle = \lambda_2 \langle v_1, v_2 \rangle \quad\Longrightarrow\quad (\lambda_1 - \lambda_2)\langle v_1, v_2 \rangle = 0
			\]
			Since $\lambda_1 \neq \lambda_2$, $\lambda_1 - \lambda_2 \neq 0$, so $\langle v_1, v_2 \rangle = 0$.
		\end{proof}
		
	\subsection{The Spectral Theorem}
	\label{sec:spectral-theorem}
		
		\begin{definition}[Invariant Subspace]
			\label{def:invariant-subspace}
			Let $A \in \mathrm{Mat}(n,\mathbb{R})$ and let $U \subseteq \mathbb{R}^n$ be a subspace. $U$ is invariant under $A$ if $Av \in U$ for every $v \in U$.
		\end{definition}
		
		\begin{definition}[Orthogonal Complement]
			\label{def:orthogonal-complement}
			Let $U \subseteq \mathbb{R}^n$ be a subspace. The orthogonal complement of $U$ is
			\[
				U^\perp := \{ w \in \mathbb{R}^n \mid \langle u, w \rangle = 0 \text{ for all } u \in U \}
			\]
		\end{definition}
		
		\begin{lemma}[Invariance of the Orthogonal Complement]
riance-complement}
			Let $A \in \mathrm{Mat}(n,\mathbb{R})$ be symmetric, and let $U \subseteq \mathbb{R}^n$ be a subspace invariant under $A$ (Definition~\ref{def:invariant-subspace}). Then $U^\perp$ (Definition~\ref{def:orthogonal-complement}) is also invariant under $A$.
		\end{lemma}
		
		\begin{proof}
			Let $A \in \mathrm{Mat}(n,\mathbb{R})$ be symmetric, and let $U \subseteq \mathbb{R}^n$ be a subspace invariant under $A$ (Definition~\ref{def:invariant-subspace}). We wish to show that $U^\perp$ (Definition~\ref{def:orthogonal-complement}) is also invariant under $A$.
			
			Let $w \in U^\perp$ be arbitrary. By Definition~\ref{def:orthogonal-complement}, our goal is to show $Aw \in U^\perp$, i.e.\
			\[
				\langle u, Aw \rangle = 0 \qquad \text{for every } u \in U
			\]
			
			Fix an arbitrary $u \in U$; we verify the equality above for this $u$.
			
			Since $A$ is symmetric, $\langle u, Av \rangle = \langle Au, v \rangle$ for all $u,v \in \mathbb{R}^n$, by Definition~\ref{def:self-adjoint} and Remark~\ref{rem:self-adjoint},
			\[
				\langle u, Aw \rangle = \langle Au, w \rangle
			\]
			
			Since $U$ is invariant under $A$ (Definition~\ref{def:invariant-subspace}) and $u \in U$,
			\[
			Au \in U
			\]
			
			We now have $\langle Au, w\rangle$ with $Au \in U$ and $w \in U^\perp$. By Definition~\ref{def:orthogonal-complement}, the inner product of any element of $U$ with any element of $U^\perp$ is zero:
			\[
				\langle Au, w \rangle = 0
			\]
			
			Chaining these equalities,
			\[
				\langle u, Aw \rangle = \langle Au, w \rangle = 0
			\]
			Since $u \in U$ was arbitrary, $Aw \in U^\perp$ by Definition~\ref{def:orthogonal-complement}. Since $w \in U^\perp$ was arbitrary, $A$ maps every vector of $U^\perp$ back into $U^\perp$; that is, $U^\perp$ is invariant under $A$ (Definition~\ref{def:invariant-subspace}).
		\end{proof}
		
		\begin{corollary}[Orthogonal Complement of the Kernel]
			\label{cor:kernel-image-orthogonal-complement}
			Let $A\in\mathrm{Mat}(m,n,\mathbb{R})$. Then
			\[
				\ker(A)^\perp = \operatorname{im}(A^\top)
			\]
		\end{corollary}
		
		\begin{proof}
			For any $w=A^\top y\in\operatorname{im}(A^\top)$ and $v\in\ker(A)$,
			\[
				\langle w,v\rangle = (A^\top y)^\top v = y^\top(Av) = y^\top 0 = 0
			\]
			so $w\in\ker(A)^\perp$. Hence $\operatorname{im}(A^\top)\subseteq\ker(A)^\perp$.
			
			By Rank-Nullity (Theorem~\ref{thm:rank-nullity}), $\dim\ker(A)=n-\operatorname{rank}(A)$. For any subspace $U\subseteq\mathbb{R}^n$, $\dim U^\perp = n-\dim U$ \citep{GolubVanLoan2013}, so $\dim\ker(A)^\perp = \operatorname{rank}(A)$. Also, row rank equals column rank, $\operatorname{rank}(A^\top) = \operatorname{rank}(A)$ by \cite{GolubVanLoan2013}, so $\dim\operatorname{im}(A^\top)=\operatorname{rank}(A)$ too. So $\ker(A)^\perp$ and $\operatorname{im}(A^\top)$ have equal dimension.
			
			Equal dimension plus containment gives equality: $\ker(A)^\perp = \operatorname{im}(A^\top)$.
		\end{proof}
		
		\begin{theorem}[Real Spectral Theorem]
			\label{thm:spectral-theorem}
			Let $A \in \mathrm{Mat}(n,\mathbb{R})$ be symmetric. Then there exists an orthonormal basis $\{q_1,\dots,q_n\}$ of $\mathbb{R}^n$ consisting of eigenvectors of $A$.
			
			Equivalently, there exists an orthogonal matrix $Q \in \mathrm{Mat}(n,\mathbb{R})$ (Definition~\ref{def:orthogonal-matrix}) with columns $q_1,\dots,q_n$, and a diagonal matrix $\Lambda \in \mathrm{Mat}(n,\mathbb{R})$ with the corresponding eigenvalues $\lambda_1,\dots,\lambda_n$ on the diagonal, such that
			\[
				A = Q\Lambda Q^\top
			\]
		\end{theorem}
		
		\begin{proof}
			We proceed by induction on $n$.
			
			\textit{Base case ($n=1$).} Let $A = [\lambda] \in \mathrm{Mat}(1,\mathbb{R})$. Any nonzero $v \in \mathbb{R}$ satisfies $Av = \lambda v$, so $v$ is an eigenvector (Definition~\ref{def:eigenvalue-eigenvector}). Setting $q_1 := v/|v|$ gives an orthonormal basis $\{q_1\}$ of $\mathbb{R}^1$. \\
			
			\textit{Inductive step.} Assume every symmetric $(n-1)\times(n-1)$ matrix has an orthonormal eigenbasis. Let $A \in \mathrm{Mat}(n,\mathbb{R})$ be symmetric. \\
			
			By Definition~\ref{def:characteristic-polynomial}, $A$'s characteristic polynomial $p_A$ has degree $n$, so by Remark~\ref{rem:existence-eigenvalues-complex} (FTA) it has $n$ roots over $\mathbb{C}$, all real by Theorem~\ref{thm:symmetric-real-eigenvalues}. Let $\lambda_1 \in \mathbb{R}$ be one such root. Since $A - \lambda_1 \mathbbm{1}$ is a real matrix with $\det(A-\lambda_1\mathbbm{1})=0$, it is not invertible, so by the Rank-Nullity Theorem (Theorem~\ref{thm:rank-nullity}) it has nontrivial kernel: there exists nonzero $v \in \mathbb{R}^n$ with $(A-\lambda_1\mathbbm{1})v = 0$, i.e.\ $Av = \lambda_1 v$. Normalising, $q_1 := v/\|v\|$ satisfies $Aq_1 = \lambda_1 q_1$ and $\|q_1\|=1$. \\
			
			Let $U := \operatorname{span}_{\mathbb{R}}\{q_1\}$; since $Aq_1=\lambda_1q_1 \in U$, $U$ is invariant under $A$ (Definition~\ref{def:invariant-subspace}), and by Lemma~\ref{lem:invariance-complement}, so is $U^\perp$ (Definition~\ref{def:orthogonal-complement}), with $\dim U^\perp = n-1$. \\
			
			By Corollary~\ref{cor:orthonormal-basis-exists}, $U^\perp$ has an orthonormal basis $\{e_1,\dots,e_{n-1}\}$ (Definition~\ref{def:orthogonal-orthonormal-set}). Since $U^\perp$ is invariant, each $Ae_j \in U^\perp$ can be expanded in this basis, defining $A' \in \mathrm{Mat}(n-1,\mathbb{R})$ by
			\[
				A'_{ij} := \langle e_i, Ae_j \rangle
			\]
			via the same coefficient-extraction argument used to construct $R$ in Theorem~\ref{thm:existence-qr-decomp} (Remark~\ref{rem:inner-product-as-matmul}). Since $A$ is symmetric, $\langle u,Av\rangle=\langle Au,v\rangle$ for all $u,v\in\mathbb{R}^n$ (Definition~\ref{def:self-adjoint}/Remark~\ref{rem:self-adjoint}); restricting to $u=e_i,v=e_j$ gives $A'_{ij} = \langle e_i,Ae_j\rangle = \langle Ae_i,e_j\rangle = \langle e_j,Ae_i\rangle = A'_{ji}$, using symmetry of the inner product \ref{ax:IP1}. So $A'$ is symmetric. \\
			
			By the inductive hypothesis, $A'$ has an orthonormal eigenbasis $\{q_2',\dots,q_n'\} \subseteq \mathbb{R}^{n-1}$. The coordinate map $\phi: \mathbb{R}^{n-1} \to U^\perp$, $\phi(v') := \sum_{i=1}^{n-1} v'_i e_i$, is an isometry, since $\{e_1,\dots,e_{n-1}\}$ is itself orthonormal: for $v',w'\in\mathbb{R}^{n-1}$, $\langle \phi(v'),\phi(w')\rangle = \sum_{i,j} v'_iw'_j\langle e_i,e_j\rangle = \sum_i v'_iw'_i = \langle v',w'\rangle$. Hence $\phi$ carries the orthonormal basis $\{q_2',\dots,q_n'\}$ to an orthonormal set $q_2,\dots,q_n \in U^\perp$, with $Aq_k=\lambda_k q_k$, $\lambda_k\in\mathbb{R}$ real by Theorem~\ref{thm:symmetric-real-eigenvalues} applied to $A'$. \\
			
			Since $q_1 \in U$ and $q_2,\dots,q_n \in U^\perp$, Definition~\ref{def:orthogonal-complement} gives $\langle q_1,q_k\rangle=0$ for $k\ge2$; together with the orthonormality of $\{q_2,\dots,q_n\}$, $\{q_1,\dots,q_n\}$ is an orthonormal set (Definition~\ref{def:orthogonal-orthonormal-set}) of $n$ eigenvectors of $A$, hence an orthonormal basis of $\mathbb{R}^n$ by Proposition~\ref{prop:orthogonal-and-linearly-independent} and Proposition~\ref{prop:v-basis}. \\
			
			Let $Q$ have columns $q_1,\dots,q_n$ (orthogonal, Definition~\ref{def:orthogonal-matrix}) and let $\Lambda := \operatorname{diag}(\lambda_1,\dots,\lambda_n)$. Then $AQ = Q\Lambda$ column-by-column, so, using $Q^{-1}=Q^\top$,
			\[
				A = Q\Lambda Q^\top
			\]
			
			This completes the induction.
		\end{proof}
		
	\begin{compremark}
		\texttt{jacobi\_eigensolver.cpp} (Appendix~\ref{app:jacobi-solver}) makes two algorithmic choices worth justifying. These are choices about how quickly and reliably the method converges in floating point, not consequences of Theorem~\ref{thm:spectral-theorem} itself. The theorem guarantees that an orthonormal eigenbasis exists, but doesn't tell us how to compute one.
		
		\begin{enumerate}[label=\roman*.]
			\item \textbf{Iteration cap scaled with $n$.} The loop stops once the largest off-diagonal entry is below tolerance, or after $50n^2$ iterations. Since each sweep scans $O(n^2)$ pairs, scaling the cap with $n^2$ keeps it proportional to the actual problem size, instead of a fixed constant that's too small for large matrices or too big for small ones.
			
			\item \textbf{Choosing the root of the same sign as $\tau$.} \texttt{computeRotationAngle()} solves for $t=\tan\theta$ using $\tau = (A_{qq}-A_{pp})/(2A_{pq})$, and picks the root with the same sign as $\tau$ (check lines 40-44). Both roots zero out $A_{pq}$ equally well algebraically, so this is a purely numerical choice: the two roots involve $\tau\pm\sqrt{1+\tau^2}$, and for large $|\tau|$, $\sqrt{1+\tau^2}\approx|\tau|$, so subtracting same-signed terms would cancel almost entirely and lose precision. Adding instead avoids this, and also picks the smaller rotation angle, which converges faster too.
		\end{enumerate}
		
		This is verified on the $5\times 5$ tridiagonal test case in \texttt{main()} by cross-checking the recovered eigenvalues printed against an independently computed trace and determinant of $A$.
	\end{compremark}
	
\section{The Singular Value Decomposition}
	
	This section develops the singular value decomposition (SVD), derived from the eigendecomposition of $A^\top A$, together with the Moore-Penrose pseudoinverse and the low-rank approximation guaranteed by the Eckart-Young theorem. \\
	
	Beyond its use elsewhere in this chapter for solving least-squares problems and characterising a matrix's best low-rank approximations, the SVD is one of the central tools of applied linear algebra: it reveals a matrix's rank, range and null space directly from its singular values and underlies dimensionality reduction techniques such as principal component analysis. \\
	
	The algorithm is implemented from scratch in C++ as \texttt{svd\_jacobi.cpp} (Appendix~\ref{app:svd-jacobi}), computing the SVD via Jacobi diagonalisation of $A^\top A$.
	
	\begin{theorem}[Singular Value Decomposition]
		\label{thm:svd}
		Let $A \in \mathrm{Mat}(m, n, \mathbb{R})$. Then there exists a singular value decomposition (SVD) i.e. an orthogonal matrix $U \in \mathrm{Mat}(m,\mathbb{R})$, an orthogonal matrix $V \in \mathrm{Mat}(n,\mathbb{R})$ and a matrix $\Sigma \in \mathrm{Mat}(m,n,\mathbb{R})$ that is diagonal with ordered entries $\sigma_1 \ge \sigma_2 \ge \cdots \ge \sigma_{\min(m,n)} \ge 0$, such that
		\[
			A = U\Sigma V^\top
		\]
		Moreover, the columns $v_1,\dots,v_n$ of $V$ are an orthonormal eigenbasis of $A^\top A$ (Theorem~\ref{thm:spectral-theorem}), the columns $u_1,\dots,u_m$ of $U$ are an orthonormal basis of $\mathbb{R}^m$, and the two are paired rather than chosen independently: for each $i$ with $\sigma_i \neq 0$,
		\[
			u_i = \frac{1}{\sigma_i} Av_i
		\]
		where $\sigma_i := \sqrt{\lambda_i}$ and $\lambda_i$ is the eigenvalue of $A^\top A$ corresponding to $v_i$.
	\end{theorem}
	
	\begin{proof}
		Consider $A^\top A \in \mathrm{Mat}(n,\mathbb{R})$. It is symmetric and $\langle A^\top A v, v\rangle = \|Av\|^2 \ge 0$ for all $v$, so all its eigenvalues are non-negative. By the Spectral Theorem (Theorem~\ref{thm:spectral-theorem}), $A^\top A$ has an orthonormal eigenbasis $v_1,\dots,v_n$ with eigenvalues $\lambda_1 \ge \cdots \ge \lambda_n \ge 0$. \\
		
		Define $\sigma_i := \sqrt{\lambda_i}$, and let $r$ be the number of strictly positive $\sigma_i$ (so $\sigma_1\ge\cdots\ge\sigma_r>0$ and $\sigma_{r+1}=\cdots=\sigma_n=0$).
		
		For $i \le r$, set $u_i := Av_i/\sigma_i$. These are orthonormal: for $j,k\le r$,
		\[
			\langle u_j,u_k\rangle = \frac{1}{\sigma_j\sigma_k}\langle Av_j,Av_k\rangle = \frac{1}{\sigma_j\sigma_k}\langle A^\top Av_j,v_k\rangle = \frac{\lambda_j}{\sigma_j\sigma_k}\langle v_j,v_k\rangle = \frac{\sigma_j}{\sigma_k}\delta_{jk}
		\]
		which is $1$ if $j=k$ and $0$ otherwise. Extend $\{u_1,\dots,u_r\}$ to a full orthonormal basis $\{u_1,\dots,u_m\}$ of $\mathbb{R}^m$. Since $\{u_1,\dots,u_r\}$ is orthonormal, it is linearly independent (Proposition~\ref{prop:orthogonal-and-linearly-independent}), so Corollary~\ref{cor:basis-extension} extends it to a basis of $\mathbb{R}^m$; applying Gram-Schmidt (Theorem~\ref{thm:gram-schmidt}) and normalising (Corollary~\ref{cor:orthonormal-basis-exists}) then gives an orthonormal basis $\{u_1,\dots,u_m\}$, with $u_1,\dots,u_r$ unchanged by Remark~\ref{rem:gram-schmidt-fixes-orthonormal}. Let $U,V$ be the matrices with these columns; both are orthogonal by construction. \\
		
		It suffices to check $AV=U\Sigma$, i.e.\ $Av_k = \sigma_k u_k$ for every $k$. For $k\le r$ this is immediate from how $u_k$ was defined. For $k>r$, $\sigma_k=0$ and $\lambda_k=0$, so $\|Av_k\|^2 = \langle A^\top Av_k,v_k\rangle = 0$, forcing $Av_k=0=\sigma_k u_k$ too. So $Av_k=\sigma_k u_k$ holds for every $k$, giving $AV=U\Sigma$. Multiplying by $V^\top=V^{-1}$ on the right gives $A=U\Sigma V^\top$.
	\end{proof}
	
	\begin{compremark}
		\texttt{svd\_jacobi.cpp}'s \texttt{computeSigmaAndU()} (Appendix~\ref{app:svd-jacobi}) uses a thin $U$ with zero columns, not the full orthonormal extension Theorem~\ref{thm:svd}'s proof builds. That proof works to extend $\{u_1,\dots,u_r\}$ into a full $m\times m$ orthogonal matrix (Corollary~\ref{cor:basis-extension}, Gram-Schmidt, Remark~\ref{rem:gram-schmidt-fixes-orthonormal}). The code skips this: \texttt{U} is $m\times n$, and columns past $r$ stay zero by default. This is safe because those columns get multiplied by zero entries of $\Sigma$ anyway, so they never affect $U\Sigma V^\top$. The full extension makes $U$ mathematically complete, but isn't needed for reconstruction.
		
		This is exercised directly in \texttt{main()}: the test matrix has rank $3$, so at least three columns of $U$ are necessarily zero, and reconstruction still matches $A$ to machine precision ($2.66\times10^{-15}$ max error).
	\end{compremark}
	
	\begin{definition}[Moore-Penrose Pseudoinverse]
		\label{def:pseudoinverse}
		Let $A \in \mathrm{Mat}(m,n,\mathbb{R})$ with SVD $A = U \Sigma V^\top$ (Theorem~\ref{thm:svd}), where $\Sigma$ has non-zero entries $\sigma_1, \dots, \sigma_r$ and zeros elsewhere. We define $\Sigma^\dagger \in \mathrm{Mat}(n, m, \mathbb{R})$ by replacing each non-zero $\sigma_i$ with the reciprocal $1/\sigma_i$ and transposing the shape. The pseudoinverse of $A$ is
		\[
			A^\dagger := V \Sigma^\dagger U^\top
		\]
	\end{definition}
	
	\begin{remark}
		If $A$ is square and invertible, then $A^\dagger = A^{-1}$. $r=n=m$, so $\Sigma^\dagger = \Sigma^{-1}$ exactly.
	\end{remark}
	
	\begin{proposition}[Least Squares Property of the Pseudoinverse]
		\label{prop:least-squares-pseudoinverse}
		Let $A \in \mathrm{Mat}(m, n, \mathbb{R})$ and $b \in \mathbb{R}^m$, and let $x^\dagger := A^\dagger b$ (Definition~\ref{def:pseudoinverse}). Then:
		\begin{enumerate}[label=\roman*.]
			\item For all $x \in \mathbb{R}^n$,
			\[
				\|Ax^\dagger - b\| \le \|Ax - b\|
			\]
			so $x^\dagger$ minimises the residual $\|Ax-b\|$ over all $x\in\mathbb{R}^n$. Equality holds if and only if $x \in x^\dagger + \ker(A)$ ($\operatorname{rank}(B)\le k$).
			\item Among all minimisers of the residual, $x^\dagger$ has the smallest norm: if $x\in\mathbb{R}^n$ satisfies $\|Ax-b\| = \|Ax^\dagger - b\|$, then
			\[
				\|x^\dagger\| \le \|x\|
			\]
			with equality if and only if $x = x^\dagger$.
		\end{enumerate}
	\end{proposition}
	
	\begin{proof}
		By Theorem~\ref{thm:svd}, $A = U\Sigma V^\top$ with $U,V$ orthogonal and $\Sigma$ having singular values $\sigma_1\ge\cdots\ge\sigma_r>0$ on the diagonal, zeros elsewhere.
		
		Let $y := V^\top x$ (so $x=Vy$, $\|x\|=\|y\|$, since $V$ is orthogonal) and $c := U^\top b$. Since $U$ is orthogonal, $\|Uz-b\|=\|z-U^\top b\|$ for any $z$; applying this with $z=\Sigma y$,
		\[
			\|Ax-b\| = \|U\Sigma V^\top x - b\| = \|\Sigma y - c\|
		\]
		Since $(\Sigma y)_i=\sigma_i y_i$ for $i\le r$ and $0$ for $i>r$,
		\[
			\|\Sigma y-c\|^2 = \sum_{i=1}^{r}(\sigma_i y_i-c_i)^2 + \sum_{i=r+1}^{m}c_i^2
		\]
		The second sum is fixed; the first depends only on $y_1,\dots,y_r$, and is minimised (to zero) exactly when $y_i=c_i/\sigma_i$ for $i\le r$. The remaining coordinates $y_{r+1},\dots,y_n$ don't affect the residual at all, so any choice of them gives a minimiser.
		
		By Definition~\ref{def:pseudoinverse}, $x^\dagger := A^\dagger b = V\Sigma^\dagger c$, so $V^\top x^\dagger = (c_1/\sigma_1,\dots,c_r/\sigma_r,0,\dots,0)$ — exactly the minimiser above, with the free coordinates set to $0$. So $x^\dagger$ minimises the residual, giving the inequality in (i).
		
		\textit{Equality condition.} $x=Vy$ achieves equality iff $y_i=c_i/\sigma_i$ for $i\le r$ and $y_{r+1},\dots,y_n$ arbitrary, i.e.\ $x = x^\dagger + \sum_{i>r}t_iv_i$ for arbitrary $t_i$.
		
		We identify this set as exactly $\ker(A)$. Since $Av_i=\sigma_iu_i=0$ for $i>r$ (shown in the proof of Theorem~\ref{thm:svd}), $\operatorname{span}_{\mathbb{R}}\{v_{r+1},\dots,v_n\}\subseteq\ker(A)$. Conversely, take any $x\in\ker(A)$ and write $x=Vy$; then
		\[
			0 = Ax = U\Sigma V^\top(Vy) = U\Sigma y
		\]
		Since $U$ is orthogonal, hence injective, this forces $\Sigma y=0$, i.e.\ $\sigma_iy_i=0$ for each $i\le r$; since $\sigma_i\neq0$ for $i\le r$, this gives $y_i=0$ for $i\le r$, with $y_{r+1},\dots,y_n$ unconstrained. So $x=Vy=\sum_{i>r}y_iv_i \in \operatorname{span}_{\mathbb{R}}\{v_{r+1},\dots,v_n\}$. Combined with the other inclusion, $\ker(A)=\operatorname{span}_{\mathbb{R}}\{v_{r+1},\dots,v_n\}$.
		
		So equality holds iff $x\in x^\dagger+\ker(A)$.
		
		\textit{Minimum norm.} If $x=x^\dagger+z$ with $z\in\ker(A)$ also minimises the residual, then $z\in\operatorname{span}_{\mathbb{R}}\{v_{r+1},\dots,v_n\}$ while $x^\dagger\in\operatorname{span}_{\mathbb{R}}\{v_1,\dots,v_r\}$, so $\langle x^\dagger,z\rangle=0$ by orthonormality of $\{v_i\}$. Then
		\[
			\|x\|^2 = \|x^\dagger+z\|^2 = \|x^\dagger\|^2+\|z\|^2 \ge \|x^\dagger\|^2,
		\]
		using bilinearity (Remark~\ref{rem:bilinearity}), with equality iff $z=0$, i.e.\ $x=x^\dagger$.
	\end{proof}
	
	\begin{lemma}[Unitary invariance of the Frobenius norm]
		\label{lem:invariance-of-frobenius}
		For any $M \in \mathrm{Mat}(m, n, \mathbb{R})$ and orthogonal matrices $U \in \mathrm{Mat}(m, \mathbb{R})$ and $V \in \mathrm{Mat}(n, \mathbb{R})$ we have:
		\[
			\|UMV^\top\|_F = \|M\|_F
		\]
		Hence multiplying $M$ by orthogonal matrices doesn't change its Frobenius norm.
	\end{lemma}
	
	\begin{proof}
		First, the trace identity $\operatorname{tr}(AB)=\operatorname{tr}(BA)$ for any $A\in\mathrm{Mat}(m,n,\mathbb{R})$, $B\in\mathrm{Mat}(n,m,\mathbb{R})$: writing out entries,
		\[
			\operatorname{tr}(AB) = \sum_i (AB)_{ii} = \sum_i\sum_j A_{ij}B_{ji} = \sum_j\sum_i B_{ji}A_{ij} = \sum_j (BA)_{jj} = \operatorname{tr}(BA)
		\]
		simply swapping the order of summation.
		
		By Definition~\ref{def:induced-norm} and Example~\ref{ex:trace-inner-product},
		\[
			\|UMV^\top\|_F^2 = \operatorname{tr}\left((UMV^\top)^\top(UMV^\top)\right)
		\]
		Since transposing a product reverses the order of factors, $(UMV^\top)^\top = VM^\top U^\top$, so
		\[
			(UMV^\top)^\top(UMV^\top) = VM^\top U^\top UMV^\top = VM^\top MV^\top
		\]
		using $U^\top U = \mathbbm{1}$, since $U$ is orthogonal (Definition~\ref{def:orthogonal-matrix}). Hence
		\[
			\|UMV^\top\|_F^2 = \operatorname{tr}(VM^\top MV^\top)
		\]
		Applying the trace identity with $X:=V$, $Y:=M^\top MV^\top$,
		\[
			\operatorname{tr}(VM^\top MV^\top) = \operatorname{tr}(XY) = \operatorname{tr}(YX) = \operatorname{tr}(M^\top MV^\top V) = \operatorname{tr}(M^\top M)
		\]
		using $V^\top V=\mathbbm{1}$, since $V$ is orthogonal, in the last step. So
		\[
			\|UMV^\top\|_F^2 = \operatorname{tr}(M^\top M) = \|M\|_F^2
		\]
		and, as both sides are non-negative, taking square roots gives $\|UMV^\top\|_F = \|M\|_F$.
	\end{proof}
	
	\begin{theorem}[Eckart-Young]
		\label{thm:eckart-young}
		Let $A \in \mathrm{Mat}(m,n,\mathbb{R})$ with SVD $A=U\Sigma V^\top$ (Theorem~\ref{thm:svd}) and singular values $\sigma_1\ge\cdots\ge\sigma_r>0$. For $0\le k<r$, we define
		\[
			A_k := \sum_{i=1}^{k} \sigma_i u_i v_i^\top
		\]
		Then for every $B \in \mathrm{Mat}(m,n,\mathbb{R})$ whose image has dimension at most $k$,
		\[
			\|A-B\|_F \ge \|A-A_k\|_F
		\]
		so $A_k$ is a best rank-$k$ approximation of $A$ in the Frobenius norm (Example~\ref{ex:trace-inner-product}). Moreover,
		\[
			\|A-A_k\|_F = \sqrt{\sigma_{k+1}^2+\cdots+\sigma_r^2}
		\]
	\end{theorem}
	
	\begin{proof}
		For any $M\in\mathrm{Mat}(m,n,\mathbb{R})$, writing its own SVD $M=U_M\Sigma_MV_M^\top$ and applying Lemma~\ref{lem:invariance-of-frobenius} gives $\|M\|_F = \|\Sigma_M\|_F = \sqrt{\sum_i\sigma_i(M)^2}$. We use this fact twice.
		
		First, $A-A_k = \sum_{i=k+1}^r\sigma_iu_iv_i^\top$ is itself an SVD, with singular values $\sigma_{k+1},\dots,\sigma_r$, so
		\[
			\|A-A_k\|_F = \sqrt{\sigma_{k+1}^2+\cdots+\sigma_r^2}
		\]
		which is also the "moreover" clause of the theorem.
		
		Now let $B$ have $\operatorname{rank}(B)\le k$. By Weyl's inequality for singular values \citep{GolubVanLoan2013}, $\sigma_{i+k}(A) \le \sigma_i(A-B) + \sigma_{k+1}(B)$ for every $i\ge1$; since $B$ has at most $k$ nonzero singular values, $\sigma_{k+1}(B)=0$, so
		\[
			\sigma_{i+k}(A) \le \sigma_i(A-B).
		\]
		Squaring and summing over $i=1,\dots,r-k$,
		\[
			\sum_{i=1}^{r-k}\sigma_{i+k}(A)^2 \le \sum_{i=1}^{r-k}\sigma_i(A-B)^2 \le \|A-B\|_F^2
		\]
		the second inequality since the right sum is only part of the full non-negative sum defining $\|A-B\|_F^2$. Reindexing $j:=i+k$, the left side is $\sum_{j=k+1}^r\sigma_j(A)^2 = \|A-A_k\|_F^2$. So
		\[
			\|A-A_k\|_F^2 \le \|A-B\|_F^2
		\]
		and taking square roots gives $\|A-A_k\|_F \le \|A-B\|_F$.
	\end{proof}
	
	